%%%%%%%%%%%%%%%%%%%%%%%%%%%%%%%%%%%%%%%%%%%%%%%%%%%%%%%%%%%%%
%%%%%%%%%%%%%%%%%%%%%%%%%%%%%%%%%%%%%%%%%%%%%%%%%%%%%%%%%%%%%
%%%%%%%%%%%%%%%%%%%%%%%%%%%%%%%%%%%%%%%%%%%%%%%%%%%%%%%%%%%%%
\part{Linear Algebra}

%%%%%%%%%%%%%%%%%%%%%%%%%%%%%%%%%%%%%%%%%%%%%%%%%%%%%%%

\chapter{Matrices and Matrix Operations in Matlab}

%%%%%%%%%%%%%%%%

You should review the vector operations in Lecture~\ref{lec:vectors}.

\subsection*{Matrix operations}

Recall how to multiply a matrix $A$ times a vector $\vb$:
$$
   A \vb = \left( \begin{array}{cc}
                    1 & 2  \\
                    3 & 4  \\
                   \end{array} \right)
                    \left( \begin{array}{c}
                    -1  \\
                    2 \\
                   \end{array} \right)
=
\left( \begin{array}{c}
                    1 \cdot (-1) + 2 \cdot 2  \\
                    3 \cdot (-1) + 4 \cdot 2 \\
                   \end{array} \right)
=
\left( \begin{array}{c}
                   3  \\
                   5 \\
                   \end{array} \right).
$$
This is a special case of matrix multiplication. To
multiply two matrices, $A$ and $B$ you proceed as follows:
$$
   A B = \left( \begin{array}{cc}
                    1 & 2  \\
                    3 & 4  \\
                   \end{array} \right)
\left( \begin{array}{cc}
                    -1 & -2 \\
                    2 &  1\\
                   \end{array} \right)
=
\left( \begin{array}{cc}
                    -1 + 4 & -2 + 2  \\
                    -3 + 8 & -6 + 4 \\
                   \end{array} \right)
=
\left( \begin{array}{cc}
                   3 & 0 \\
                   5 & -2 \\
                   \end{array} \right).
$$
Here both $A$ and $B$ are $2 \times 2$ matrices.
Matrices can be multiplied together in this way provided
that the number of columns of $A$ match the number of rows
of $B$. We always list the size of a matrix by rows, then
columns, so a $3 \times 5$ matrix would have 3 rows and 5 columns.
So, if $A$ is $m \times n$ and $B$ is $p \times q$, then
we can multiply $AB$ if and only if $n = p$.
A column vector can be thought of as a $p \times 1$ matrix
and a row vector as a $1 \times q$ matrix. Unless otherwise specified
we will assume a vector $\vb$ to be a column vector and
so $A\vb $ makes sense as long as the number of columns of $A$
matches the number of entries in $\vb$.

Printing matrices on the screen takes up a lot of space, so you may
want to use
\begin{lstlisting}[frame=none,numbers=left]
format compact
\end{lstlisting}
Enter a matrix into Matlab either as
\begin{lstlisting}[frame=none,numbers=left]
A = [ 1  3 -2 5 ;  -1  -1 5 4 ; 0 1 -9  0]
\end{lstlisting}
or
\begin{lstlisting}[frame=none,numbers=left]
A = [1,3,-2,5; -1,-1,5,4; 0,1,-9,0]
\end{lstlisting}
Also enter a vector $\ub$:
\begin{lstlisting}[frame=none,numbers=left]
u = [ 1  2  3  4]'
\end{lstlisting}
To multiply a matrix times a vector $A\mathbf{u}$ use \verb&*&:
\begin{lstlisting}[frame=none,numbers=left]
A*u
\end{lstlisting}
Since $A$ is 3 by 4 and $\ub$ is 4 by 1 this multiplication is
valid and the result is a 3 by 1 vector.

Now enter another matrix $B$ using
\begin{lstlisting}[frame=none,numbers=left]
B = [3 2 1; 7 6 5; 4 3 2]
\end{lstlisting}
You can multiply $B$ times $A$ with
\begin{lstlisting}[frame=none,numbers=left]
B*A
\end{lstlisting}
but $A$ times $B$ is not defined and
\begin{lstlisting}[frame=none,numbers=left]
A*B
\end{lstlisting}
will result in an error message.


You can multiply a matrix by a scalar:
\begin{lstlisting}[frame=none,numbers=left]
2*A
\end{lstlisting}
Adding matrices $A + A$ will give the same result:
\begin{lstlisting}[frame=none,numbers=left]
A + A
\end{lstlisting}
You can even add a number to  a matrix:
\begin{lstlisting}[frame=none,numbers=left]
A + 3    % add 3 to every entry of A
\end{lstlisting}

%%%%%%%%%%%%%%%%

\subsection*{Component-wise operations}

Just as for vectors, adding a `\verb&.&' before `\verb&*&',
`\verb&/&', or `\verb&^&' produces entry-wise multiplication,
division and exponentiation.
If you enter
\begin{lstlisting}[frame=none,numbers=left]
B*B
\end{lstlisting}
the result will be $BB$, i.e.\ matrix multiplication of $B$ times itself. But,
if you enter
\begin{lstlisting}[frame=none,numbers=left]
B.*B
\end{lstlisting}
the entries of the resulting matrix will contain the squares
of the same entries of $B$. Similarly if you want $B$  multiplied
by itself 3 times then enter
\begin{lstlisting}[frame=none,numbers=left]
B^3
\end{lstlisting}
but, if you want to cube all the entries of B then enter
\begin{lstlisting}[frame=none,numbers=left]
B.^3
\end{lstlisting}
Note that \verb&B*B& and \verb&B^3& only make sense if $B$ is
square, but \verb&B.*B& and \verb&B.^3& make sense for any size
matrix.


%%%%%%%%%%%%%%%%

\subsection*{The identity matrix and the inverse
of a matrix}

The $n \times n$ {\em identity matrix} is a square matrix with
ones on the diagonal and zeros everywhere else. It is
called the identity because it plays the same role that $1$
plays in multiplication, i.e.
$$
      A I = A,   \qquad I A = A, \qquad I \vb = \vb
$$
for any matrix $A$ or vector $\vb$ where the sizes match.
An identity matrix in \textsc{Matlab} is produced by the command
\begin{lstlisting}[frame=none,numbers=left]
I = eye(3)
\end{lstlisting}

A square matrix $A$ can have an {\em inverse} which is denoted
by $A^{-1}$. The definition of the inverse is that
$$
       A A^{-1} = I  \qquad \text{and} \qquad A^{-1} A = I.
$$
In theory an inverse is very important, because if you
have an equation
$$
     A \xb = \bb
$$
where $A$ and $\bb$ are known and $\xb$ is unknown (as
we will see, such problems are very common and important)
then the theoretical solution is
$$
      \xb = A^{-1} \bb.
$$
We will see later that this is not a practical way to solve
an equation, and $A^{-1}$ is only important for the purpose
of derivations.
%
In \textsc{Matlab} we can calculate a matrix's inverse
very conveniently:
\begin{lstlisting}[frame=none,numbers=left]
C = randn(5,5)
inv(C)
\end{lstlisting}
%
However, not all square matrices have inverses:
\begin{lstlisting}[frame=none,numbers=left]
D = ones(5,5)
inv(D)
\end{lstlisting}


%%%%%%%%%%%%%%%%


\subsection*{The ``Norm'' of a matrix}

For a vector, the ``norm'' means the same thing as the length (geometrically, not the number of entries).
Another way to think of it is how far the vector is from being
the zero vector.
We want to measure a matrix in much the same way and the {\em
norm} is such a quantity. The usual definition of the norm of
a matrix is
\begin{defn}
Suppose $A$ is a $m \times n$ matrix. The norm of $A$ is
$$
  \|A\| \equiv  \max_{\|\vb\|=1} \|A\vb\|.
$$
\end{defn}
The maximum in the definition is taken over all vectors with
length $1$ (unit vectors), so the definition means the largest
factor that the matrix stretches (or shrinks) a unit vector.
This definition seems cumbersome at first, but it
turns out to be the best one. For example, with this definition
we have the following inequality for any vector $\vb$:
$$
    \|A\vb\| \le \|A\| \|\vb\|.
$$
In \textsc{Matlab} the norm of a matrix is obtained by the
command
\begin{lstlisting}[frame=none,numbers=left]
norm(A)
\end{lstlisting}
For instance the norm of an identity matrix is $1$:
\begin{lstlisting}[frame=none,numbers=left]
norm(eye(100))
\end{lstlisting}
and the norm of a zero matrix is $0$:
\begin{lstlisting}[frame=none,numbers=left]
norm(zeros(50,50))
\end{lstlisting}

For a matrix the norm defined above and calculated by
\textsc{Matlab} is not the square root of the sum
of the square of its entries. That quantity is called the
{\em Froebenius norm}, which is also sometimes useful,
but we will not need it.

%%%%%%%%%%%%%%%%%

\subsection*{Some other useful commands}

% Try out the following:
% \begin{lstlisting}[frame=none,numbers=left]
% C = rand(5,5)    % random matrix with uniform distribution in [0,1]
% size(C)          % gives the dimensions (m x n) of C
% det(C)           % the determinant of the matrix
% max(C)           % the maximum of each column
% min(C)           % the minimum in each column
% sum(C)           % sums each column
% mean(C)          % the average of each column
% diag(C)          % just the diagonal elements
% C'               % tranpose the matrix
% \end{lstlisting}
\verb&C = rand(5,5)& \dotfill random matrix with uniform
         distribution in $[0,1]$.\\
\verb&size(C)& \dotfill gives the dimensions ($m\times n$) of $C$.\\
\verb&det(C)&   \dotfill the determinant of the matrix.\\
\verb&max(C)&   \dotfill the maximum of each column.\\
\verb&min(C)&   \dotfill the minimum in each column.\\
\verb&sum(C)&   \dotfill sums each column.\\
\verb&mean(C)&   \dotfill the average of each column.\\
\verb&diag(C)&   \dotfill just the diagonal elements.\\
\verb&C'&   \dotfill tranpose the matrix.

In addition to \verb&ones&, \verb&eye&, \verb&zeros&, \verb&rand& and
\verb&randn&,  \textsc{Matlab} has several other commands that
automatically produce special matrices:
%\begin{lstlisting}[frame=none,numbers=left]
%hilb(6)
%pascal(5)
%\end{lstlisting}
\\
\verb&hilb(6)&\\
\verb&pascal(5)&


%%%%%%%%%%%%
\newpage
\subsection*{Exercises}
\begin{list}
{\arabic{chapter}.\arabic{exercise}}{\usecounter{exercise}}


\item Enter the matrix $A$ by
\begin{lstlisting}[frame=none,numbers=left]
A = [3 2 1; 6 5 4; 9 8 7]
\end{lstlisting}
and also the matrix
$$
B = \left[
{\begin{array}{rrr}
7 & 8 & 9 \\
4 & 5 & 6 \\
1 & 2 & 3 
\end{array}}
\right]\,.  $$
Find  (a) \verb&A*A&, (b) \verb&A^2&, (c) \verb&A.^2&, (d) \verb&A.*B&, (e) \verb&A*B&. Turn in the output.

\item By hand, calculate $A \vb$, $AB$, and $BA$ for:
      $$
        A = \left( \begin{array}{ccc}
                    2 & 4 & -1 \\
                    -2 & 1 & 9 \\
                    -1 & -1 & 0 \\
                   \end{array} \right),
        \quad
        B = \left( \begin{array}{ccc}
                    0 & -1 & -1 \\
                    1 & 0 & 2 \\
                    -1 & -2 & 0 \\
                   \end{array} \right),
        \quad
        \vb = \left( \begin{array}{c}
                    3  \\
                    1  \\
                    -1 \\
                   \end{array} \right).
     $$
     Check the results using \textsc{Matlab}. Think about how
     fast computers are. Turn in your hand work.

\item  Write a well-commented \textsc{Matlab} {\bf function} program
    \verb&myinvcheck& that
    \begin{itemize}
    \item makes a $n \times n$ random matrix (normally distributed,
      \verb&A = randn(n,n)&),
    \item calculates its inverse (\verb&B = inv(A)&),
    \item multiplies the two back together,
    \item calculates the residual (difference between $AB$ and \verb&eye(n)&), and
    \item returns the scalar residual (norm of the difference).
    \end{itemize}
    Turn in your program.
  
 \item Write a well-commented \textsc{Matlab} {\bf script} program \verb&myinvcheckplot& that calls
    \verb&myinvcheck& for $n=10,20,40,\ldots,2^i10$ for $i$ as large as possible, 
    records the
    results of each trial, and plots the scalar residual versus \verb&n& 
    using a log plot. (See \verb&help loglog&.)
 
   What happens to the scalar residual as \(n\)
  gets big? Turn in the program, the plot, and a very brief report on
  the results of your experiments.  (Do not print any large random matrices.) 

\end{list}


%%%%%%%%%%%%%%%%%%%%%%%%%%%%%%%%%%%%%%%%%%

\chapter{Introduction to Linear Systems}
\label{chp:intro_linsys}
%%%%%%%%%%%%%%%%

\subsection*{How linear systems occur}


Linear systems of equations naturally occur in many places
in engineering, such as structural analysis, dynamics and
electric circuits. Computers have made it possible to
quickly and accurately solve larger and larger systems
of equations. Not only has this allowed engineers to handle
more and more complex problems where linear systems naturally
occur, but has also prompted engineers to use linear systems
to solve problems where they do not naturally occur such
as thermodynamics, internal stress-strain analysis, fluids and
chemical processes. It has become standard practice in
many areas to analyze a problem by transforming it into a linear
systems of equations and then solving those equation by computer.
In this way, computers have made linear systems of equations
the most frequently used tool in modern engineering.


In Figure~\ref{equilateral_truss} we show a truss with equilateral
triangles. In Statics you may use the ``method of joints'' to write
equations for each node of the truss\footnote{See
  http://en.wikipedia.org/wiki/Truss or
  http://en.wikibooks.org/wiki/Statics for reference.}.  This set of
equations is an example of a linear system.  Making the approximation
$\sqrt{3}/2 \approx .8660$, the equations for this truss are
\begin{equation}
\label{truss_equilateral_eqns}
\begin{split}
  .5 \, T_1 + T_2 & = R_1 = f_1 \\
  .866 \, T_1 & = - R_2 = -.433 \, f_1 - .5 \, f_2 \\
  -.5 \, T_1 + .5 \, T_3 + T_4 &= - f_1 \\
  .866 \, T_1 + .866 \, T_3 &= 0 \\
  -T_2 - .5 \, T_3 + .5 \, T_5 + T_6 &= 0 \\
  .866 \, T_3 + .866 \, T_5 &= f_2 \\
  -T_4 - .5 \, T_5 + .5 \, T_7 &=0, \\
\end{split}
\end{equation}
where $T_i$ represents the tension in the $i$-th member of the truss.
%
% \begin{figure}
% \setlength{\unitlength}{1 mm}
% \LARGE
% \begin{center}
% \begin{picture}(64,17)(0,-1)
% 	\put(0,0){\line(1,0){64}}
% 	\put(16,16){\line(1,0){32}}
% 	\put(0,0){\line(1,1){16}}
% 	\put(16,0){\line(1,1){16}}
% 	\put(64,0){\line(-1,1){16}}
% 	\put(48,0){\line(-1,1){16}}
% 	\put(16,0){\line(0,1){16}}
% 	\put(48,0){\line(0,1){16}}
% 	\put(0,0){\circle*{2}}
% 	\put(16,0){\circle*{1}}
% 	\put(48,0){\circle*{1}}
% 	\put(64,0){\circle*{2}}
% 	\put(16,16){\circle*{1}}
% 	\put(32,16){\circle*{1}}
% 	\put(48,16){\circle*{1}}
% \end{picture}
% \end{center}
% \caption{A truss.}
% \label{truss}
% \end{figure}



\begin{figure}[ht]
\setlength{\unitlength}{1 mm}
% \LARGE
 \begin{center}
 \begin{picture}(64,40)(0,-10)
% nodes
 	\put(0,0){\circle*{2}}
 	\put(32,0){\circle*{1}}
 	\put(64,0){\circle*{2}}
 	\put(16,27.7){\circle*{1}}
 	\put(48,27.7){\circle*{1}}
	\put(3,1){A}
	\put(14,21){B}
	\put(31,4){C}
	\put(47,22){D}
	\put(59,1){E}
% edges
 	\put(0,0){\line(1,0){64}}
 	\put(16,27.7){\line(1,0){32}}
	\qbezier(0,0)(8,13.85)(16,27.7)
	\qbezier(32,0)(40,13.85)(48,27.7)
	\qbezier(16,27.7)(24,13.85)(32,0)
	\qbezier(48,27.7)(56,13.85)(64,0)
	\put(4,12){1}
	\put(15,1){2}
	\put(26,12){3}
	\put(31,24){4}
	\put(41,12){5}
	\put(47,1){6}
	\put(57,13){7}
% forces
 	\put(6,27.7){\vector(1,0){6}}
	\put(2,27){$f_1$}
 	\put(32,-1){\vector(0,-1){6}}
	\put(31,-10){$f_2$}
 	\put(0,-7){\vector(0,1){5}}
	\put(-1,-10){$R_2$}
 	\put(64,-7){\vector(0,1){5}}
	\put(62,-10){$R_3$}
 	\put(-2,0){\vector(-1,0){6}}
	\put(-13,-1){$R_1$}
% 	\put(16,0){\line(1,1){16}}
% 	\put(64,0){\line(-1,1){16}}
% 	\put(48,0){\line(-1,1){16}}
% 	\put(16,0){\line(0,1){16}}
% 	\put(48,0){\line(0,1){16}}
 \end{picture}
\end{center}
%\centerline{\hbox{\psfig{figure=equilateral_truss3.eps,height=1.5in,width=2.5in}}}
\caption{An equilateral truss. Joints or nodes are labeled alphabetically, $A$, $B$, \ldots
and Members (edges) are labeled numerically: $1$, $2$, \ldots. The forces $f_1$ and $f_2$
are applied loads and $R_1$, $R_2$ and $R_3$ are reaction forces applied by the supports.}
\label{equilateral_truss}
\end{figure}



You could solve
this system by hand with a little time and patience; systematically
eliminating variables and substituting.
Obviously, it would be a lot better to put the equations
on a computer and let the computer solve it. In the next few lectures
we will learn how to use a computer effectively to solve linear systems.
The first key to dealing with linear systems is to realize that
they are equivalent to matrices, which contain numbers, not variables.

As we discuss various aspects of matrices, we wish to keep in
mind that the matrices that come up in engineering systems are
{\em really  large}. It is not unusual in real engineering
to use matrices whose dimensions are in the thousands!
It is frequently the case that a method that is fine for
a $2 \times 2$ or $3 \times 3$ matrix is totally inappropriate
for a $2000 \times 2000$ matrix. We thus want to emphasize methods
that work for large matrices.


%%%%%%%%%%%%%%%

\subsection*{Linear systems are equivalent to matrix equations}
\label{lec:LSME}

The system of linear equations
\begin{equation*}
\begin{split}
     x_1 - 2 x_2 + 3 x_3 & = 4 \\
     2 x_1 - 5x_2 + 12 x_3 & = 15 \\
           2 x_2 -  10 x_3 & = -10 
\end{split}
\end{equation*}
is equivalent to the matrix equation
$$
\left(\begin{array}{ccc}
        1 & -2 & 3  \\
        2 & -5 & 12 \\
        0 & 2 & -10
      \end{array}\right)
\left(\begin{array}{c}
       x_1 \\ x_2 \\ x_3
      \end{array} \right)
=
\left(\begin{array}{c}
       4\\ 15 \\ -10
      \end{array} \right),
$$
which is equivalent to the {\em augmented matrix}
$$
\left(\begin{array}{ccc|c}
        1 & -2 & 3 & 4 \\
        2 & -5 & 12 & 15 \\
        0 & 2 & -10 & -10
      \end{array}\right).
$$
The advantage of the augmented matrix, is that it contains only
numbers, not variables. The reason this is better is because computers
are much better in dealing with numbers than variables. To solve
this system, the main steps are called {\em Gaussian elimination} and
{\em back substitution}.

The augmented matrix for the equilateral truss equations
(\ref{truss_equilateral_eqns}) is given by
\begin{equation}\label{truss_equilateral_matrix}
\left(\begin{array}{ccccccc|c}
        .5    & 1 & 0 & 0 & 0 & 0 & 0 & f_1 \\
        .866  & 0 & 0  & 0 & 0 & 0 &  0 & -.433  f_1 - .5 \, f_2  \\
         -.5  & 0 & .5 &  1 & 0 & 0 & 0  & - f_1 \\
        .866  & 0 & .866  & 0 & 0 & 0 & 0 & 0 \\
         0    & -1 & - .5  & 0 & .5 & 1 &  0 & 0 \\
         0    & 0  & .866  & 0 & .866 & 0 & 0 &  f_2 \\
         0 &  0 &  0  & -1  & - .5  &0 & .5 & 0 \\
      \end{array}\right).
\end{equation}
Notice that a lot of the entries are $0$. Matrices like this, called
{\em sparse}, are common in applications and there are methods
specifically designed to efficiently handle sparse matrices.


%%%%%%%%%%%%%%%%

\subsection*{Triangular matrices and back substitution}

Consider a linear system whose augmented matrix happens to be
\begin{equation}\label{upper}
\left(\begin{array}{ccc|c}
        1 & -2 & 3 & 4 \\
        0 & -1 & 6 & 7 \\
        0 & 0 & 2 & 4
      \end{array}\right).
\end{equation}
Recall that each row represents an equation and each column a variable.
The last row represents the equation $2 x_3 = 4$. The equation is
easily solved, i.e. $x_3 = 2$. The second row represents the equation
$- x_2 + 6 x_3 = 7$, but since we know $x_3 = 2$, this simplifies to:
$ - x_2 + 12 = 7$. This is easily solved, giving $x_2 = 5$. Finally,
since we know $x_2$ and $x_3$, the first row simplifies to:
$x_1 - 10 + 6 = 4$. Thus we have $x_1 = 8$ and so we know the
whole solution vector: $\xb = \langle 8 , 5, 2 \rangle$. The process
we just did is called {\em back substitution}, which is both efficient
and easily programmed. The property that made it possible
to solve the system so easily is that $A$ in this case is {\em upper
triangular}. In the next section we show an efficient way to
transform an augmented matrix into an upper triangular matrix.


%%%%%%%%%%%%%%%%

\subsection*{Gaussian Elimination}

Consider the matrix
$$
 A = \left(\begin{array}{ccc|c}
        1 & -2 & 3  &  4 \\
        2 & -5 & 12 & 15 \\
        0 & 2 & -10 & -10
      \end{array}\right).
$$
The first step of Gaussian elimination is to get rid of the 2 in
the (2,1) position by subtracting $2$ times the first row
from the second row, i.e. {\sf (new 2nd = old 2nd - (2) 1st)}.
We can do this because it is essentially the
same as adding equations, which is a valid
algebraic operation.  This leads to
$$
\left(\begin{array}{ccc|c}
        1 & -2 & 3 & 4 \\
        0 & -1 & 6 & 7 \\
        0 & 2 & -10 & -10
      \end{array}\right).
$$
There is already a zero in the lower left corner, so we don't need to
eliminate anything there. To eliminate the third row, second column,
we need to subtract $-2$ times the second row from the third row, {\sf
  (new 3rd = old 3rd - (-2) 2nd)}, to obtain
$$
\left(\begin{array}{ccc|c}
        1 & -2 & 3 &  4 \\
        0 & -1 & 6 &  7 \\
        0 &  0 & 2 &  4
      \end{array}\right).
$$
This is now just exactly the matrix in equation (\ref{upper}), which
we can now solve by back substitution.


%%%%%%%%%%%%%%%%

\subsection*{\textsc{Matlab}'s matrix solve command}

In \textsc{Matlab} the standard way to solve a system $A \xb = \bb$
is by the command
\begin{lstlisting}[frame=none,numbers=left]
x = A \ b
\end{lstlisting}
This syntax is meant to suggest dividing by \(A\) from the left as in
\[A \xb = \bb \quad \Leftrightarrow\quad A\backslash A \xb = A\backslash \bb 
\quad \Leftrightarrow\quad  \xb = A\backslash \bb
\,.\]
Such division is not meaningful mathematically, but it helps for remembering the syntax.

This command carries out Gaussian elimination and back substitution.
We can do the above computations as follows:
\begin{lstlisting}[frame=none,numbers=left]
A = [1 -2 3 ; 2 -5 12 ; 0 2 -10]
b = [4 15 -10]'
x = A \ b
\end{lstlisting}

Next, use the \textsc{Matlab} commands above to solve $A \xb = \bb$ when
the augmented matrix for the system is
 $$
\left(\begin{array}{ccc|c}
        1 & 2 & 3 &  4 \\
        5 & 6 & 7 &  8 \\
        9 &  10 & 11 &  12
      \end{array}\right),
$$
by entering
\begin{lstlisting}[frame=none,numbers=left]
x1 = A \ b
\end{lstlisting}
Check the result by entering
\begin{lstlisting}[frame=none,numbers=left]
A*x1 - b
\end{lstlisting}
You will see that the resulting answer satisfies the equation exactly.
Next try solving using the inverse of $A$:
\begin{lstlisting}[frame=none,numbers=left]
x2 = inv(A)*b
\end{lstlisting}
This answer can be seen to be inaccurate by checking
\begin{lstlisting}[frame=none,numbers=left]
A*x2 - b
\end{lstlisting}
Thus we see one of the reasons why the inverse is never used for actual
computations, only for theory.

%%%%%%%%%%%%%%%%

\subsection*{Exercises}
\nobreak
\begin{list}
{\arabic{chapter}.\arabic{exercise}}{\usecounter{exercise}}

\item  Set $f_1 = 1000N$ and $f_2 = 5000N$ in
the equations (\ref{truss_equilateral_eqns}) for the equailateral truss.
Input the coefficient matrix $A$ and the right hand side vector $b$ in
(\ref{truss_equilateral_matrix}) into \textsc{Matlab}.
Solve the system using the command \verb&\& to find the tension in each member of the
truss. Save the matrix $A$ as \verb&A_equil_truss& and keep it for later use.
(Enter \verb& save A_equil_truss A&.)
Print out and turn in $A$, $\bb$ and the solution $\xb$.
\label{exr:truss}

\item Write each system of equations as an augmented matrix. By hand, find
the solutions using {\em Gaussian elimination and back substitution} (i.e.\ the exact
algorithm in this chapter). Check
your solutions using \textsc{Matlab}.

\begin{enumerate}[(a)]
\item 
\begin{equation*}
\begin{split}
     x_1 +  x_2   & = 2 \\
   4 x_1 + 5 x_2  & = 10
\end{split}
\end{equation*}
\item
\begin{equation*}
\begin{split}
     x_1 + 2 x_2 + 3 x_3 & = -1 \\
   4 x_1 + 7 x_2 + 14 x_3 & = 3 \\
    x_1 +  4 x_2 + 4 x_3 & =  1
\end{split}
\end{equation*}
\end{enumerate}

\end{list}



%%%%%%%%%%%%%%%%%%%%%%%%%%%%%%%%%%%%%%%%%%%

\chapter{Some Facts About Linear Systems}
\label{sec:lineartheory}


%%%%%%%%%%%%%%%%

\subsection*{Some inconvenient truths}


In the last lecture we learned how to solve a linear system
using \text{Matlab}.
Input the following:
\begin{lstlisting}[frame=none,numbers=left]
A = ones(4,4)
b = randn(4,1)
x = A \ b
\end{lstlisting}

As you will find, there is no solution to the equation $A \xb = \bb$.
This unfortunate circumstance is mostly the fault of the matrix, $A$,
which is too simple, its columns (and rows) are all the same.
Now try
\begin{lstlisting}[frame=none,numbers=left]
b = ones(4,1)
x = [ 1 0 0 0]'
A*x
\end{lstlisting}
So the system $A \xb = \bb$ does have a solution. Still unfortunately,
that is not the only solution. Try
\begin{lstlisting}[frame=none,numbers=left]
x = [ 0 1 0 0]'
A*x
\end{lstlisting}
We see that this $x$ is also a solution. Next try
\begin{lstlisting}[frame=none,numbers=left]
x = [ -4  5  2.27 -2.27]'
A*x
\end{lstlisting}
This $x$ is a solution! It is not hard to see that
there are endless possibilities for solutions of this
equation.

%%%%%%%%%%%%%%%

\subsection*{Basic theory}


The most basic theoretical fact about linear systems is
\begin{thm}
A linear system $A\xb = \bb$ may have 0, 1, or infinitely many solutions.
\end{thm}
In most (but not all!) engineering applications we would want to
have exactly one solution. The following two theorems tell us
exactly when we can and cannot expect this.
\begin{thm}
Suppose $A$ is a square ($n \times n $) matrix.
The following are all equivalent:
\begin{enumerate}
\item The equation $A\xb = \bb$ has exactly one solution for any $\bb$.
\item $det(A) \ne 0$.
\item $A$ has an inverse.
\item The only solution of $A\xb = \mathbf{0}$ is $\xb = \mathbf{0}$.
\item The columns of $A$ are linearly independent (as vectors).
\item The rows of $A$ are linearly independent.
\end{enumerate}
If $A$ has these properties then it is called {\bf non-singular}.
\end{thm}

On the other hand, a matrix that does not have these properties
is called {\bf singular}.
\begin{thm}\label{def:singular}
Suppose $A$ is a square  matrix.
The following are all equivalent:
\begin{enumerate}
\item The equation $A\xb = \bb$ has 0 or $\infty$ many solutions
  depending on $\bb$.
\item $\det(A) = 0$.
\item $A$ does not have an inverse.
\item The equation $A\xb = \mathbf{0}$ has solutions other than $\xb =
  \mathbf{0}$.
\item The columns of $A$ are linearly dependent as vectors.
\item The rows of $A$ are linearly dependent.
\end{enumerate}
\end{thm}


To see how the two theorems work, define two matrices (type in
\verb&A1& then scroll
up and modify to make \verb&A2&) :
$$
\verb&A1& = \left(\begin{array}{ccc}
        1 & 2 & 3  \\
        4 & 5 & 6 \\
        7 & 8 & 9
      \end{array}\right),
\qquad
\verb&A2& = \left(\begin{array}{ccc}
        1 & 2 & 3  \\
        4 & 5 & 6 \\
        7 & 8 & 8
      \end{array}\right),
$$
and two vectors:
$$
\verb&b1&  = \left(\begin{array}{c}
        0 \\ 3 \\ 6
             \end{array}\right),
\qquad
\verb&b2&  = \left(\begin{array}{c}
        1 \\ 3 \\ 6
             \end{array}\right).
$$
First calculate the determinants of the matrices:
\begin{lstlisting}[frame=none,numbers=left]
det(A1)
det(A2)
\end{lstlisting}
Then attempt to find the inverses:
\begin{lstlisting}[frame=none,numbers=left]
inv(A1)
inv(A2)
\end{lstlisting}
Which matrix is singular and which is non-singular?
Finally, attempt to solve all the possible equations $A \xb = \bb$:
\begin{lstlisting}[frame=none,numbers=left]
x = A1 \ b1
x = A1 \ b2
x = A2 \ b1
x = A2 \ b2
\end{lstlisting}
As you can see, equations involving the non-singular matrix have one
and only one solution, but equation involving a singular matrix
are more complicated.

%%%%%%%%%%%%%%%%

\subsection*{The Residual}

Recall that the residual for an approximate solution $x$
of an equation $f(x) =0$ is defined as $r = \|f(x)\|$.
It is a measure of how close the equation is to being satisfied.
For a linear system of equations we define the residual vector
of an approximate solution $\xb$ by
\begin{equation*}
%\label{lsres}
\rb = A \xb - \bb.
\end{equation*}
If the solution vector $\xb$ were exactly correct, then $\rb$ would be 
exactly the zero vector.  The size (norm)  of $\rb$ is
an indication of how close we have come to solving
$A \xb = \bb$. We will refer to this number as the {\bf scalar residual} or just  the {\bf residual} of the
approximation solution:
\begin{equation}
   r = \| A \xb - \bb \|.
\end{equation}


%%%%%%%%%%%%%%%

\subsection*{Exercises}
\nobreak
\begin{list}
  {\arabic{chapter}.\arabic{exercise}}{\usecounter{exercise}}


\item By hand, find all the solutions (if any) of the following linear
  system using the augmented matrix and Gaussian elimination (following exactly
  the algorithm in the notes):
  \begin{equation*}
    \begin{split}
      x_1 + 2 x_2 + 3 x_3 & = 4, \\
      4 x_1 + 5 x_2 + 6 x_3 & = 10, \\
      7 x_1 + 8 x_2 + 9 x_3 & = 14 \,.
    \end{split}
  \end{equation*}
Try solving this system in \textsc{Matlab} using the command \verb&x = A \ b&. What happens?
Turn in your hand work.

\item
  \begin{enumerate}
  \item[(a)] Write a well-commented \textsc{Matlab} {\bf function}
    program \verb&mysolvecheck& with input a number \verb&n& that
    makes a random $n \times n$ matrix $A$ and a random vector $\bb$,
    solves the linear system $A \xb = \bb$, calculates the scalar
    residual $r = \|A \xb - \bb\|$, and outputs that number as  $r$.
  \item[(b)] Write a well-commented \textsc{Matlab} {\bf script}
    program that calls \verb&mysolvecheck& 10 times each for $n = $
    5, 10, 20, 40, 80, and 160, then records and averages the
    results and makes a log-log plot of the average $r$ vs.\ $n$. 
    Once your program is running correctly, increase the maximum $n$ (by factors of $2$) until the program stops running 
    within 5 minutes.
  \end{enumerate}
  Turn in the plot and the two programs. (Do not print any large random matrices.)
\end{list}


%%%%%%%%%%%%%%%%%%%%%%%%%%%%%%%%%%%%%%%%%%%%%%%%

\chapter{Accuracy, Condition Numbers and Pivoting}

In this lecture we will discuss two separate issues
of accuracy in solving linear systems. The first,
{\em pivoting}, is a method that ensures that Gaussian
elimination proceeds as accurately as possible.
The second, {\em condition number}, is a measure
of how bad a matrix is. We will see that if a matrix has
a bad condition number, the solutions are unstable with
respect to small changes in data.

%%%%%%%%%%%%%

\subsection*{The effect of rounding}

All computers store numbers as  finite strings of binary
floating point digits (bits). This limits numbers to a fixed
number of significant digits and implies that after even
the most basic calculations, rounding must happen.

Consider the following exaggerated example. Suppose
that our computer can only store 2 significant digits
and it is asked to do Gaussian elimination on
\begin{equation*}
\left(\begin{array}{cc|c}
        .001 & 1 & 3  \\
          1 &  2 & 5
      \end{array}\right).
\end{equation*}
Doing the elimination exactly would produce
\begin{equation*}
\left(\begin{array}{cc|c}
        .001 & 1 & 3  \\
          0 & -998 & -2995
      \end{array}\right),
\end{equation*}
but rounding to 2 digits, our computer would store this as
\begin{equation*}
\left(\begin{array}{cc|c}
        .001 & 1 & 3  \\
          0 &  -1000 & -3000
      \end{array}\right).
\end{equation*}
Backsolving this reduced system gives
$$
     x_1 = 0 \quad\text{and}\quad x_2 = 3.
$$
This seems fine until you realize that backsolving the unrounded
system gives
$$
    x_1 = -1 \quad\text{and}\quad x_2 = 3.001.
$$

%%%%%%%%%%%%%%%

\subsection*{Row Pivoting}


A way to fix the problem is to use pivoting, which means to switch rows of
the matrix. Since switching rows of the augmented matrix
just corresponds to switching the order of the equations, no harm
is done:
\begin{equation*}
\left(\begin{array}{cc|c}
          1 &  2 & 5  \\
        .001 & 1 & 3
      \end{array}\right).
\end{equation*}
Exact elimination would produce
\begin{equation*}
\left(\begin{array}{cc|c}
          1 &  2 & 5  \\
          0 & .998 & 2.995
      \end{array}\right).
\end{equation*}
Storing this result with only 2 significant digits gives
\begin{equation*}
\left(\begin{array}{cc|c}
          1 &  2 & 5  \\
          0 & 1 & 3
      \end{array}\right).
\end{equation*}
Now backsolving produces
$$
      x_1 = -1 \quad\text{and}\quad x_2 = 3,
$$
which is the true solution (rounded to 2 significant digits).

The reason this worked is because $1$ is bigger than $0.001$. To pivot
we switch rows so that the largest entry in a column is the one used
to eliminate the others. In bigger matrices, after each column is completed,
compare the diagonal element of the next column with all the entries below
it. Switch it (and the entire row) with the one with greatest absolute value.
For example in the following matrix, the first column is finished and before
doing the second column, pivoting should occur since $|-2|>|1|$:
$$
\left(\begin{array}{ccc|c}
        1 & -2 & 3 & 4 \\
        0 &  1 & 6 & 7 \\
        0 & -2 & -10 & -10
      \end{array}\right).
$$
Pivoting the 2nd and 3rd rows would produce
$$
\left(\begin{array}{ccc|c}
        1 & -2 & 3 & 4 \\
        0 & -2 & -10 & -10\\
        0 &  1 & 6 & 7
      \end{array}\right).
$$

%%%%%%%%%%

\subsection*{Condition number}

In some systems, problems occur even without rounding.
Consider the following augmented matrices:
$$
\left(\begin{array}{cc|c}
        1   & 1/2 & 3/2  \\
        1/2 & 1/3 &  1  \\
      \end{array}\right)
     \quad \text{and} \quad
\left(\begin{array}{cc|c}
        1   & 1/2  & 3/2 \\
        1/2 & 1/3  & 5/6\\
      \end{array}\right).
$$
Here we have the same $A$, but two different input vectors:
$$  \bb_1 = (3/2,1)' \quad \textrm{and} \quad \bb_2 = (3/2,5/6)'
$$
which are pretty close to one another. You would
expect then that the solutions $\xb_1$ and $\xb_2$ would also be close.
Notice that this matrix does not need pivoting. Eliminating exactly
we get
$$
\left(\begin{array}{cc|c}
        1 & 1/2  & 3/2  \\
        0 & 1/12 & 1/4 \\
      \end{array}\right)
     \quad \text{and} \quad
\left(\begin{array}{cc|c}
        1 & 1/2  & 3/2 \\
        0 & 1/12 & 1/12\\
      \end{array}\right).
$$
Now solving we find
$$
   \xb_1 = (0,3)' \quad \textrm{and} \quad \xb_2 = (1,1)'
$$
which are {\em not close at all} despite the fact that we did the
calculations exactly. This poses a new problem: some matrices
are very sensitive to small changes in input data. The extent of
this sensitivity is measured by the {\bf condition number}.
The definition of condition number is: consider all small changes
$\delta A$ and $\delta \bb$ in $A$ and $\bb$  and the resulting
change, $\delta \xb$, in the solution $\xb$. Then
$$
      \textrm{cond}(A) \equiv \max \left( \frac{ \|\delta \xb\|/\|\xb\|}
         { \frac{\|\delta A\|}{\|A\|} + \frac{\|\delta \bb\|}{\|\bb\|} } \right) =
     \max \left( \frac{\text{Relative error of output}}
                      {\text{Relative error of inputs}} \right).
$$
Put another way, changes in the input data get multiplied by the condition number
to produce changes in the outputs.
Thus a high condition number is bad. It implies that small errors
in the input can cause large errors  in the output.
%
It is not obvious from our definition above, but one can prove that the condition number of a matrix is at least 1.


In \textsc{Matlab} enter
\begin{lstlisting}[frame=none,numbers=left]
H = hilb(2)
\end{lstlisting}
which should result in the matrix above. \textsc{Matlab} produces the
condition number of a matrix with the command
\begin{lstlisting}[frame=none,numbers=left]
cond(H)
\end{lstlisting}
Thus for this matrix small errors in the input can get magnified
by 19 in the output.
Next try the matrix
\begin{lstlisting}[frame=none,numbers=left]
A = [ 1.2969 0.8648 ; .2161  .1441]
cond(A)
\end{lstlisting}
For this matrix small errors in the input can get magnified
by $2.5 \times 10^8$ in the output! (We will see this happen in the exercise.)
This is obviously not very good for engineering
where all measurements, constants and inputs are approximate.

Is there a solution to the problem of bad condition numbers?
Usually, bad condition numbers in engineering contexts result
from poor design. So, the engineering solution to bad conditioning
is {\bf redesign}.

Finally, find the determinant of the matrix $A$ above:
\begin{lstlisting}[frame=none,numbers=left]
det(A)
\end{lstlisting}
which will be small. If $\det(A) = 0$ then the matrix is singular,
which is bad because it implies there will not be a unique solution.
The case here, $\det(A) \approx 0$, is also bad, because it means the
matrix is almost singular. Although $\det(A) \approx 0$ generally
indicates that the condition number will be large, they are actually
independent things. To see this, find the determinant and condition
number of the matrix \verb&[1e-10,0;0,1e-10]& and
the matrix \verb&[1e+10,0;0,1e-10]&.

%%%%%%%%%%%%%%%%%%%%%%%%
\newpage

\subsection*{Exercises}
\begin{list}
{\arabic{chapter}.\arabic{exercise}}{\usecounter{exercise}}

%%%
\item Let $$
A = \left[
{\begin{array}{cc}
1.2969 & .8648 \\
.2161 & .1441
\end{array}}
 \right]\,.
$$
\begin{enumerate}
\item[(a)] Find the condition number, determinant and inverse of $A$ (using \textsc{Matlab}).
\item[(b)] Let $B$ be the matrix obtained from $A$ by rounding off to three
  decimal places ($1.2969\mapsto 1.297$).  Find the inverse of $B$.  
  How do $A^{-1}$ and $B^{-1}$ differ? Explain how
  this happened.
\item[(c)] Set \verb&b1 = [1.2969; 0.2161] & and do \verb&x = A \ b1 &.
  Repeat the process but with a vector \verb&b2& obtained from
  \verb&b1& by rounding off to three decimal places. Explain exactly
  what happened.  Why was the first answer so simple?  Why do the two
  answers differ by so much?
\end{enumerate}

%%%%%%%%%%
\item To see how to solve linear systems symbolically, try
\begin{lstlisting}[frame=none,numbers=left]
B = [sin(sym(1)) sin(sym(2)); sin(sym(3)) sin(sym(4))]
c = [1; 2]
x = B \ c
pretty(x)
\end{lstlisting}

Now input the matrix
\[
\textsf{Cs = }\left[
\begin{array}{rr}
1 & 2 \\
2 & 4
\end{array}
\right]
\]
symbolically as above by wrapping each number in \verb&sym&. Create a numerical version via
\verb&Cn = double(Cs)& and define the two vectors \verb$ d1 = [4; 8] $ and
\verb$ d2 = [1; 1]$.
Solve the systems (as in the third line of code above) \verb&Cs*x = d1&, \verb&Cn*x = d1&,
\verb&Cs*x = d2&, and \verb&Cn*x = d2&.
Explain the results.  Does the symbolic or
non-symbolic way give more information?
%
%\item Recall the matrix $A$ that you saved using \verb&A_equil_truss&
%  in exercise \ref{lec:LSME}.\ref{exr:truss}.  (Enter
%  \verb&> load A_equil_truss& or click on the file
%  \verb&A_equil_truss.mat& in the folder where you saved it; this
%  should reproduce the matrix $A$.)

%  Find the condition number and determinant of this matrix.  
  
%  Write a well-commented
%  {\bf script} program to calculate the mean and median condition
%  number of 1000 randomly generated 7 by 7 matrices.  How does the
%  condition number of the truss matrix compare to randomly generated
%  matrices? Turn in your script with the answers to these questions.

%\textsf{The matrix in 1.\ is nearly singular, causing the linear system to be
%very sensitive to perturbations.  Students are exposed to both
%symbolic and numerical solutions.  The ideas of no solutions or
%infinitely many solutions are reinforced.}

\end{list}

%%%%%%%%%%%%%%%%%%%%%%%%%%%%%%%%%%%%%%%%%%%%%%%%%%%%%%%%%%%%%%

\chapter{LU Decomposition}

In many applications where linear systems appear, one needs to solve
$A \xb = \bb$ for many different vectors $\bb$. For instance, a
structure must be tested under several different loads, not just one.
As in the example of a truss (\ref{truss_equilateral_matrix}),
the loading in such a problem is usually
represented by the vector $\bb$.  Gaussian elimination with pivoting
is the most efficient and accurate way to solve a linear system.  Most
of the work in this method is spent on the matrix $A$ itself.  If we
need to solve several different systems with the same $A$, and $A$ is
big, then we would like to avoid repeating the steps of Gaussian
elimination on $A$ for every different $\bb$. This can be accomplished
by the {\em LU decomposition}, which in effect records the steps of
Gaussian elimination.

%%%%%%%%%%%%%%%%

\subsection*{LU decomposition}

The main idea of the LU decomposition
is to record the steps used in Gaussian elimination on A
in the places where the zero is produced.
Consider the matrix
$$
 A = \left(\begin{array}{ccc}
        1 & -2 & 3 \\
        2 & -5 & 12  \\
        0 & 2 & -10
      \end{array}\right).
$$
The first step of Gaussian elimination is to subtract $2$ times the
first row from the second row. In order to record what we have done,
we will put the multiplier, $2$, into the place it was used to make a
zero, i.e. the second row, first column. In order to make it clear
that it is a record of the step and not an element of $A$, we will put
it in parentheses. This leads to
$$
\left(\begin{array}{ccc}
        1 & -2 & 3 \\
        (2) & -1 & 6  \\
        0 & 2 & -10
      \end{array}\right).
$$
There is already a zero in the lower left corner, so we don't need to eliminate
anything there. We record this fact with a $(0)$. To eliminate the third row,
second column, we need to subtract
$-2$ times the second row from the third row. Recording the $-2$ in the spot
it was used we have
$$
\left(\begin{array}{ccc}
        1 & -2 & 3 \\
        (2) & -1 & 6  \\
        (0) & (-2) & 2
      \end{array}\right).
$$
Let $U$ be the upper triangular matrix produced, and let $L$ be the lower
triangular matrix with the records and ones on the diagonal, i.e.
$$
L = \left(\begin{array}{ccc}
        1 & 0 & 0 \\
        2 & 1 & 0  \\
        0 & -2 & 1
      \end{array}\right)
\quad\text{and}\quad
U = \left(\begin{array}{ccc}
         1 & -2 & 3 \\
         0 & -1 & 6  \\
         0 &  0 & 2
      \end{array}\right),
$$
then we have the following wonderful property:
$$
  LU = \left(\begin{array}{ccc}
        1 & 0 & 0 \\
        2 & 1 & 0  \\
        0 & -2 & 1
      \end{array}\right)
      \left(\begin{array}{ccc}
        1 & -2 & 3 \\
        0 & -1 & 6  \\
        0 & 0 & 2
      \end{array}\right)
=   \left(\begin{array}{ccc}
        1 & -2 & 3 \\
        2 & -5 & 12  \\
        0 & 2 & -10
      \end{array}\right)
= A.
$$
Thus we see that $A$ is actually the product of $L$ and $U$.
Here $L$ is lower triangular and $U$ is upper triangular.
When a matrix can be written as a product of simpler
matrices, we call that a {\em decomposition} of $A$ and
this one we call the LU decomposition.

%%%%%%%%%%%%%%%%%%

\subsection*{Using LU to solve equations}

If we also include pivoting, then an LU decomposition for $A$
consists of three matrices $P$, $L$ and $U$ such that
\begin{equation}\label{LUP}
   PA = LU.
\end{equation}
The pivot matrix $P$ is the identity matrix, with the same rows switched
as the rows of $A$ are switched in the pivoting. For instance,
$$
P = \left(\begin{array}{ccc}
        1 & 0 & 0 \\
        0 & 0 & 1  \\
        0 & 1 & 0
      \end{array}\right),
$$
would be the pivot matrix if the second and third rows of $A$ are
switched by pivoting.  \textsc{Matlab} will produce an $LU$
decomposition with pivoting for a matrix $A$
with the command\\
\verb& > [L U P] = lu(A)&\\
where \verb&P& is the pivot matrix. To use this information to solve
$A \xb = \bb$ we first pivot both sides by multiplying by the pivot
matrix:
$$
     PA \xb = P \bb \equiv \mathbf{d}.
$$
Substituting $LU$ for $PA$ we get
$$
     LU \xb = \mathbf{d} .
$$
Then we define $\yb=U \xb$, which is unknown since \(\xb\) is unknown. 
Using forward-substitution, we can (easily) solve
$$
     L \yb = \mathbf{d}
$$
for \(\yb\) and then using back-substitution we can (easily) solve 
$$
     U \xb = \yb
$$
for \(\xb\).
In \textsc{Matlab} this would work as follows:
\begin{lstlisting}[frame=none,numbers=left]
A = rand(5,5)
[L U P] = lu(A)
b = rand(5,1)
d = P*b
y = L\d
x = U\y
rnorm = norm(A*x - b)  % Check the result
\end{lstlisting}
We can then solve for any other $\bb$ without redoing the LU step.
Repeat the sequence for a new right hand side: \verb&c = randn(5,1)&;
you can start at the third line.
While this may not seem like a big savings, it would be if $A$ were a large
matrix from an actual application.


The LU decomposition is an example of {\em Matrix Decomposition} which
means taking a general matrix $A$ and breaking it down into components
with simpler properties.  Here $L$ and $U$ are simpler because they
are lower and upper triangular. There are many other matrix
decompositions that are useful in various contexts. Some of the most
useful of these are the QR decomposition, the Singular Value
decomposition and Cholesky decomposition. Often a decomposition
is associated with an algorithm, e.g., finding the LU decomposition is equivalent
to completing Gaussian Elimination.


%%%%%%%%%%%%%%

\subsection*{Exercises}
\begin{list}
{\arabic{chapter}.\arabic{exercise}}{\usecounter{exercise}}

\item Solve the systems below by hand using Gaussian elimination  and
back substitution (exactly as above) on the augmented matrix. 
As a by-product, give the LU decomposition
of A. Pivot wherever appropriate (the number being eliminated should be smaller than the 
number eliminating it). Check by hand that
$LU=PA$ and $A\xb = \bb$ and compare with \textsc{Matlab}.
\begin{enumerate}
\item[(a)] $\displaystyle{ A =
      \left(\begin{array}{cc} 0.5 & 4 \\
                               2 & 4
           \end{array} \right), \qquad \bb = \left(\begin{array}{c} 0 \\
                                                                -3 \end{array} \right)}
$
\item[(b)] $\displaystyle{ A =
      \left(\begin{array}{ccc} 1 & 2 & 3 \\
                             4 & 1 & 1 \\
                             2 & 3 & 9
           \end{array} \right), \qquad \bb = \left(\begin{array}{c} 2 \\
                                                                                                1 \\
                                                                                                1  \end{array} \right)}
$
\end{enumerate}


\item
Finish the following \textsc{Matlab} function program:
\begin{lstlisting}
function [x1, r1, x2, r2] = mysolve(A,b)
    % Solves linear systems using the LU decomposition with pivoting
    % and also with the built-in solve function A\b.
    % Inputs: A -- the matrix
    %         b -- the right-hand vector
    % Outputs: x1 -- the solution using the LU method
    %          r1 -- the scalar residual using the LU method
    %          x2 -- the solution using the built-in method
    %          r2 -- the scalar residual using the 
    %                built-in method
\end{lstlisting}

Using \verb&format long&, test the program on both random matrices
(\verb&randn(n,n)&) and Hilbert matrices (\verb&hilb(n)&) with $n$
large (as big as you can make it and the program still run). Print
your program and summarize your observations.  (Do not print any
random matrices or vectors.)
\end{list}


%%%%%%%%%%%%%%%%%%%%%%%%%%%%%%%%%%%%%%%%%%%%%%%%%

\chapter{Nonlinear Systems - Newton's Method}
\label{lec:multnewton}

%%%%%%%%%%%%%%%%

\subsection*{An Example}

The LORAN (LOng RAnge Navigation) system calculates the position
of a boat at sea using signals from fixed transmitters. From
the time differences of the incoming signals, the boat obtains
differences of distances to the transmitters. This leads
to two equations each representing hyperbolas defined
by the differences of distance of two points (foci). An example
of such equations\footnote{E.~Johnston, J.~Mathews, {\em Calculus}, Addison-Wesley, 2001} are
\begin{equation}\label{eqn:loran}
\begin{split}
   \frac{x^2}{186^2} - \frac{y^2}{300^2 - 186^2} &= 1
\quad\text{and}\\
  \frac{(y-500)^2}{279^2} - \frac{(x-300)^2}{500^2 - 279^2} &= 1
\,.
\end{split}
\end{equation}


Solving two quadratic equations with two unknowns would require
solving a 4 degree polynomial equation. We could do this by hand, but
for a navigational system to work well, it must do the calculations
automatically and numerically. We note that the Global Positioning
System (GPS) works on similar principles and must do similar
computations.

%%%%%%%%%%%%

\subsection*{Vector Notation}

In general, we can usually find solutions to a system of equations
when the number of unknowns matches the number of equations.
Thus, we wish to find solutions to systems that have the form
\begin{equation}\label{nlsys}
\begin{split}
f_1(x_1, x_2, x_3, \ldots , x_n) & = 0 \\
f_2(x_1, x_2, x_3, \ldots , x_n) & = 0 \\
f_3(x_1, x_2, x_3, \ldots , x_n) & = 0 \\
      \vdots \hspace*{1.4cm}      &  \\
f_n(x_1, x_2, x_3, \ldots , x_n) & = 0. \\
\end{split}
\end{equation}

For convenience we can think of $(x_1, x_2, x_3, \ldots , x_n)$ as
a vector $\xb$ and $(f_1, f_2, \ldots, f_n)$ as a vector-valued function
$\fb$. With this notation, we can write the system of equations (\ref{nlsys})
simply as
$$
   \fb(\xb) = \mathbf{0},
$$
i.e.\ we wish to find a vector that makes the vector function equal to
the zero vector.

As in Newton's method for one variable, we need to start with an
initial guess $\xb_0$. In theory, the more variables one has, the
harder it is to find a good initial guess. In practice, this must
be overcome by using physically reasonable assumptions about
the possible values of a solution, i.e.~take advantage of engineering
knowledge of the problem.

%%%%%%%%%%%%

\subsection*{Linear Approximation for Vector Functions}

In the single variable case, Newton's method was derived
by considering the linear approximation of the function $f$ at
the initial guess $\xb_0$.
From Calculus, the following is the linear approximation of
$\fb$ at $\xb_0$, for vectors and vector-valued functions:
$$
\fb(\xb) \approx \fb(\xb_0) + D\fb(\xb_0) (\xb - \xb_0).
$$
Here $D\fb(\xb_0)$ is an $n \times n$ matrix whose entries are the
various partial derivative of the components of $\fb$, evaluated at
\(\xb_0\). Specifically,
\begin{equation}
  \label{eq:Df}
   D\fb(\xb_0) = \left( \begin{array}{ccccc}
     \frac{\partial f_1}{\partial x_1}(\xb_0) & \frac{\partial f_1}{\partial x_2}(\xb_0) &
              \frac{\partial f_1}{\partial x_3}(\xb_0) & \ldots & \frac{\partial f_1}{\partial x_n}(\xb_0) \vspace*{.3cm} \\

     \frac{\partial f_2}{\partial x_1}(\xb_0) & \frac{\partial f_2}{\partial x_2}(\xb_0) &
              \frac{\partial f_2}{\partial x_3}(\xb_0) & \ldots & \frac{\partial f_2}{\partial x_n}(\xb_0) \vspace*{.15cm}\\
        \vdots & \vdots  & \vdots & \ddots & \vdots \vspace*{.15cm}\\
     \frac{\partial f_n}{\partial x_1}(\xb_0) & \frac{\partial f_n}{\partial x_2}(\xb_0) &
              \frac{\partial f_n}{\partial x_3}(\xb_0) & \ldots & \frac{\partial f_n}{\partial x_n}(\xb_0)
                       \end{array} \right).  
\end{equation}

%%%%%%%%%%%%

\subsection*{Newton's Method}

We wish to find $\xb$ that makes $\fb$ equal to the zero vectors, so
let's choose $\xb_1$ so that
$$
  \fb(\xb_0) + D\fb(\xb_0) (\xb_1 -\xb_0) = \mathbf{0}.
$$
Since $D\fb(\xb_0)$ is a square matrix, we can solve this equation
by
$$
\xb_1 = \xb_0 - (D\fb(\xb_0))^{-1} \fb(\xb_0),
$$
provided that the inverse exists. The formula is the vector equivalent
of the Newton's method formula we learned before.
However, in practice we never use the inverse of a matrix for
computations, so we cannot use this formula directly. Rather,
we can do the following. First solve the equation
\begin{equation}
D\fb(\xb_0) \Delta \xb = -\fb(\xb_0)\,,
\end{equation}
where we want to have
$$
\Delta \xb = \xb_1 - \xb_0.
$$
Since $D\fb(\xb_0)$ is a known matrix and $-\fb(\xb_0)$ is a known
vector, this equation is just a system of linear equations,
which can be solved efficiently and accurately. Once we have the
solution vector $\Delta \xb$, we can obtain our improved estimate
$\xb_1$ by
$$
   \xb_1 = \xb_0 + \Delta \xb.
$$

For subsequent steps, we have the following process:
\begin{itemize}
\item Solve $D\fb(\xb_i) \Delta \xb = -\fb(\xb_i)$ for $\Delta \xb$.
\item Let $\xb_{i+1} = \xb_i + \Delta \xb$
\end{itemize}

%%%%%%%%%%%%%%

\subsection*{An Experiment}

We will solve the following set of equations:
\begin{equation}
\begin{split}
x^3 + y &= 1 \\
y^3 - x & = -1.
\end{split}
\end{equation}
You can easily check that $(x,y) = (1,0)$ is a solution of this system.
By graphing both of the equations you can also see that $(1,0)$ is the
only solution (Figure~\ref{fig:nlsys}).
\begin{figure}
\centerline{\includegraphics[width=0.6\textwidth]{nlsys}}
\caption{Graphs of the equations $x^3 + y = 1$ and  $y^3 - x  = -1$.
There is one and only one intersection; at $(x,y) = (1,0)$.}
\label{fig:nlsys}
\end{figure}

We can put these equations into vector-function form (\ref{nlsys})
by letting $x_1 = x$,
$x_2 = y$ and
\begin{equation*}
\begin{split}
f_1(x_1,x_2) &= x_1^3 + x_2 - 1 \\
f_2(x_1,x_2) &= x_2^3 - x_1 + 1\,,
\end{split}
\end{equation*}
or, equivalently,
\begin{equation*}
\fb(\xb) = \left(\begin{array}{c}
                   x_1^3 + x_2 - 1 \\
                   x_2^3 - x_1 + 1
                  \end{array} \right).
\end{equation*}
The partial derivatives are
\begin{equation}
\label{eqn:partials}
\frac{\partial f_1}{\partial x_1} = 3 x_1^2, 
\quad \frac{\partial f_1}{\partial x_2} = 1, 
\quad \frac{\partial f_2}{\partial x_1} = -1,  
\quad\text{and}\quad \frac{\partial f_2}{\partial x_2} = 3 x_2^2.
\end{equation}
Thus,
\begin{equation}
\label{eqn:Dfexample}
D \mathbf{f} =  \left(\begin{array}{cc}
                   3x_1^2 & 1 \\
                    -1 & 3x_2^2 
                  \end{array} \right).
\end{equation}
Now that we have the equations in  vector-function form,
we can write the following script program:
\begin{lstlisting}
format long;
f = @(x)[ x(1)^3+x(2)-1 ; x(2)^3-x(1)+1 ]
x = [.5;.5]
x = fsolve(f,x)
\end{lstlisting}
Save this program as \verb&myfsolve.m& and run it.
You will see that the internal \textsc{Matlab} solving
command \verb&fsolve& approximates the solution, but
only to about 7 decimal places. While that would be close enough
for most applications, one would expect that we could do
better on such a simple problem.

Next we will implement Newton's method for this problem.
Modify your \verb&myfsolve& program to:
\begin{lstlisting}
% mymultnewton - solve an example nonlinear system
format long;
n=8  % set some number of iterations, may need adjusting
f = @(x)[x(1)^3+x(2)-1 ; x(2)^3-x(1)+1] % the vector function
% the matrix of partial derivatives
Df = @(x)[3*x(1)^2,  1 ; -1,  3*x(2)^2] 
x = [.5;.5]             % starting guess
for i = 1:n
      Dx = -Df(x)\f(x); % solve for increment
      x = x + Dx        % add on to get new guess
      f(x)              % see if f(x) is really zero
end
\end{lstlisting}

Save and run this program (as \verb&mymultnewton&) and you will see that it
finds the root exactly (to machine precision)
in only 6 iterations. Why is this simple program able to do better
than \textsc{Matlab}'s built-in program?

%%%%%%%%%%%

\subsection*{Exercises}
\begin{list}
{\arabic{chapter}.\arabic{exercise}}{\usecounter{exercise}}

\item 
  \begin{enumerate}[(a)]
  \item Put the LORAN equations (\ref{eqn:loran}) into the
    function form (\ref{nlsys}).
  \item  Calculate (by hand) the partial derivatives of $\mathbf{f}$
  and construct the matrix $D\fb$  (see (\ref{eqn:partials}) and (\ref{eqn:Dfexample})).
  \item Adapt the \verb&mymultnewton& program to find a solution
    for these equations.  By trying different starting vectors, find
    at least three different solutions. (There are actually four
    solutions.)  
   \item Think of at least one way that the navigational
    system could determine the correct solution.
  \end{enumerate}

\end{list}







%%%%%%%%%%%%%%%%%%%%%%%%%%%%%%%%%%%%%%%%%%%%%%%%%

\chapter{Eigenvalues and Eigenvectors}

Suppose that $A$ is a square ($n \times n$) matrix.
We say that a  nonzero vector $\vb$ is an eigenvector
and a number $\lambda$ is its eigenvalue if
\begin{equation}\label{def:eig}
   A \vb = \lambda \vb.
\end{equation}

Geometrically this means that $A\vb$ is in the same direction
as $\vb$, since multiplying a vector by a number changes
its length, but not its direction.\\

\textsc{Matlab} has a built-in routine for finding eigenvalues and eigenvectors:
\begin{lstlisting}[frame=none,numbers=left]
A = pascal(4)
[v e] = eig(A)
\end{lstlisting}
The results are a matrix \verb&v& that contains eigenvectors as
columns and a diagonal matrix \verb&e& that contains eigenvalues on
the diagonal. We can check this by
\begin{lstlisting}[frame=none,numbers=left]
v1 = v(:,1)
A*v1
e(1,1)*v1
\end{lstlisting}

%%%%%%%%%%%%%%%%

\subsection*{Finding Eigenvalues for $2\times 2$ and $3\times3$}


If $A$ is $2\times 2$ or $3\times3$ then we can find its eigenvalues
and eigenvectors by hand. Notice that Equation (\ref{def:eig})
can be rewritten as
\begin{equation*}
     A \vb - \lambda \vb = \mathbf{0}.
\end{equation*}
It would be nice to factor out the $\vb$ from the right-hand side of this
equation, but we can't because $A$ is a matrix and $\lambda$ is a number.
However, since $I \vb = \vb$, we can do the following:
\begin{equation*}
\begin{split}
A\vb - \lambda \vb &=    A \vb - \lambda I \vb \\
                   &= (A  - \lambda I) \vb \\
                   &= \mathbf{0}
\end{split}
\end{equation*}
If $\vb$ is nonzero, then by Theorem~\ref{def:singular} in
Lecture~\ref{sec:lineartheory} the matrix $(A - \lambda I)$ must be
singular. By the same theorem, we must have
\begin{equation*}%\label{charac}
         \det(A - \lambda I) = 0.
\end{equation*}
This is called the {\em characteristic equation}.

For a $2 \times 2$ matrix, $A - \lambda I$ is calculated as in the following example:
\begin{equation*}
\begin{split}
 A - \lambda I &=  \left(\begin{array}{cc} 1 & 4 \\
                             3 & 5
                         \end{array} \right)
                  - \lambda  \left(\begin{array}{cc} 1 & 0 \\
                                                    0 & 1
                                   \end{array} \right) \\
               &=  \left(\begin{array}{cc} 1 & 4 \\
                                           3 & 5
                         \end{array} \right)
                  - \left(\begin{array}{cc} \lambda & 0 \\
                                           0 & \lambda
                         \end{array} \right) \\
               &= \left(\begin{array}{cc} 1 -\lambda & 4 \\
                                          3 & 5 - \lambda
                        \end{array} \right)
\,.
\end{split}
\end{equation*}
The determinant of $A - \lambda I$ is then
\begin{equation*}
\begin{split}
\det(A - \lambda I) &= (1 - \lambda)(5 - \lambda) - 4 \cdot 3 \\
                   &= -7 - 6 \lambda + \lambda^2 .
\end{split}
\end{equation*}
The characteristic equation $\det(A - \lambda I) = 0$ is simply a quadratic
equation:
\begin{equation*}
           \lambda^2 - 6 \lambda - 7 = 0.
\end{equation*}
The roots of this equation are $\lambda_1 = 7$ and $\lambda_2 = -1$.
These are the eigenvalues of the matrix $A$. Now to find the corresponding
eigenvectors we return to the equation $(A - \lambda I) \vb = \mathbf{0}$.
For $\lambda_1 = 7$, the equation for the eigenvector
$(A - \lambda I) \vb = \mathbf{0}$ is equivalent
to the augmented matrix
\begin{equation}\label{eveqn}
\left(\begin{array}{cc|c} -6 & 4 & 0\\
                             3 & -2 & 0
                    \end{array} \right) .
\end{equation}
Notice that the first and second rows of this matrix are multiples of
one another. Thus Gaussian elimination would produce all zeros on the bottom row.
Thus this equation has infinitely many solutions, i.e.~infinitely
many eigenvectors. Since only the direction of the eigenvector matters, this is okay,
we only need to find one of the eigenvectors. Since the second row of the augmented
matrix represents the equation
$$
               3 x - 2 y = 0,
$$
we can let
$$
     \vb_1 = \left(\begin{array}{c} 2 \\
                             3
                    \end{array} \right).
$$
This comes from noticing that $(x,y) = (2,3)$ is a solution of
$ 3x -2y =0 $.

For $\lambda_2 = -1$, $(A - \lambda I) \vb = \mathbf{0}$ is equivalent
to the augmented matrix
$$
\left(\begin{array}{cc|c} 2 & 4 & 0\\
                             3 & 6 & 0
                    \end{array} \right)
.
$$
Once again the first and second rows of this matrix are multiples of
one another. For simplicity we can
let
$$
     \vb_2 = \left(\begin{array}{c} -2 \\
                             1
                    \end{array} \right).
$$

One can always check an eigenvector and eigenvalue by multiplying:
\begin{align*}
A \vb_1 &= \left(\begin{array}{cc} 1 & 4 \\
                             3 & 5
                    \end{array} \right)
          \left(\begin{array}{c} 2 \\
                                 3
                    \end{array} \right)
         = \left(\begin{array}{c} 14 \\
                                   21
                    \end{array} \right)
         = 7 \left(\begin{array}{c} 2 \\
                                     3
                    \end{array} \right)
	=7 \vb_1
\quad\text{and}\\
A \vb_2 &= \left(\begin{array}{cc} 1 & 4 \\
                             3 & 5
                    \end{array} \right)
          \left(\begin{array}{c} -2 \\
                                 1
                    \end{array} \right)
         = \left(\begin{array}{c} 2 \\
                                   -1
                    \end{array} \right)
         = -1 \left(\begin{array}{c} -2 \\
                                      1
                    \end{array} \right)
	= -1 \vb_2
\,.
\end{align*}

For a $3 \times 3$ matrix we could complete the same process.
The $\det(A -\lambda I)=0$ would be a cubic polynomial and we
would expect to usually get 3 roots, which are the eigenvalues.


%%%%%%%%%%%%%%%%

\subsection*{Larger Matrices}

For a $n \times n$ matrix with $n \ge 4$ this process is too
long and cumbersome to complete by hand. Further, this process
is not well suited even to implementation on a computer program
since it involves determinants and solving a $n$-degree polynomial.
For $n \ge 4$ we need more ingenious methods. These methods rely
on the geometric meaning of eigenvectors and eigenvalues rather than solving
algebraic equations. We will overview these methods in Lecture~\ref{EWmeth}.



%%%%%%%%%%%%%%%%

\subsection*{Complex Eigenvalues}

It turns out that the eigenvalues of some matrices are complex
numbers, even when the matrix only contains real numbers.
When this happens the complex eigenvalues must occur
in conjugate pairs, i.e.
$$
    \lambda_{1,2} = \alpha \pm i \beta.
$$
The corresponding eigenvectors must also come
in conjugate pairs:
$$
     \wb = \ub \pm i \vb.
$$
In applications, the imaginary part of the eigenvalue, $\beta$,
often is related to the frequency of an oscillation.
This is because of Euler's formula
$$
    e^{\alpha + i \beta} = e^\alpha(\cos \beta + i \sin \beta).
$$

Certain kinds of matrices that arise in applications can only
have real eigenvalues and eigenvectors. The most common such type of matrix
is the symmetric matrix. A matrix is symmetric if it is equal
to its own transpose, i.e.~it is symmetric across the diagonal.
For example,
$$
\left(\begin{array}{cc} 1 & 3 \\
                        3 & -5
                    \end{array} \right)
$$
is symmetric and so we know beforehand that its eigenvalues will be real,
not complex.


%%%%%%%%%%%%%%%%%%%

\subsection*{Exercises}
\begin{list}
{\arabic{chapter}.\arabic{exercise}}{\usecounter{exercise}}

\item Find the eigenvalues and eigenvectors of the following
      matrix by hand:
$$
A = \left(\begin{array}{cc} 2 & 1 \\
                            1 & 2
                    \end{array} \right).
$$


\item Find the eigenvalues and eigenvectors of the following
      matrix by hand:
$$
B = \left(\begin{array}{cc} 1 & -3 \\
                            3 &  1
                    \end{array} \right).
$$
Can you guess the eigenvalues of the matrix
$$
C = \left(\begin{array}{cc} a & -b \\
                            b &  a
                    \end{array} \right)?
$$
\end{list}



%%%%%%%%%%%%%%%%%%%%%%%%%%%%%%%%%%%%%%%%%%%%%%%%%%%%%%%%%%%

\chapter[Vibrational Modes and Frequencies]{An Application of Eigenvectors: Vibrational Modes and Frequencies}

One application of eigenvalues and eigenvectors is in the analysis of vibration
problems. A simple nontrivial vibration problem is the motion of two
objects with equal masses $m$ attached to each other and fixed outer
walls by equal springs with spring constants $k$, as shown in
Figure~\ref{fig:twosprings}.
\begin{figure}
\setlength{\unitlength}{1 mm}
%\LARGE
\begin{center}
\begin{picture}(64,10)(0,-5)
	%left wall
	\put(0,-5){\line(1,0){2}}
	\put(0,5){\line(1,0){2}}
	\put(2,-5){\line(0,1){10}}
	% spring
	\put(2,0){\line(1,0){3}}
	\qbezier(5,0)(9,5)(9,-1)
	\qbezier(9,-1)(9,-3)(8,-3)
	\qbezier(8,-3)(7,-3)(8,0)
	\qbezier(8,0)(12,5)(12,-1)
	\qbezier(12,-1)(12,-3)(11,-3)
	\qbezier(11,-3)(10,-3)(11,0)
	\qbezier(11,0)(15,5)(15,-1)
	\qbezier(15,-1)(15,-3)(14,-3)
	\qbezier(14,-3)(13,-3)(14,0)
	\qbezier(14,0)(18,5)(18,-1)
	\qbezier(18,-1)(18,-3)(17,-3)
	\qbezier(17,-3)(16,-3)(17,0)
	\qbezier(17,0)(19,4)(19,0)
	% weight
	\put(19,0){\line(1,0){6}}
	\put(22,0){\circle*{2}}
	% spring
	\qbezier(25,0)(29,5)(29,-1)
	\qbezier(29,-1)(29,-3)(28,-3)
	\qbezier(28,-3)(27,-3)(28,0)
	\qbezier(28,0)(32,5)(32,-1)
	\qbezier(32,-1)(32,-3)(31,-3)
	\qbezier(31,-3)(30,-3)(31,0)
	\qbezier(31,0)(35,5)(35,-1)
	\qbezier(35,-1)(35,-3)(34,-3)
	\qbezier(34,-3)(33,-3)(34,0)
	\qbezier(34,0)(38,5)(38,-1)
	\qbezier(38,-1)(38,-3)(37,-3)
	\qbezier(37,-3)(36,-3)(37,0)
	\qbezier(37,0)(39,4)(39,0)
	%weight
	\put(39,0){\line(1,0){6}}
	\put(42,0){\circle*{2}}
	% spring
	\qbezier(45,0)(49,5)(49,-1)
	\qbezier(49,-1)(49,-3)(48,-3)
	\qbezier(48,-3)(47,-3)(48,0)
	\qbezier(48,0)(52,5)(52,-1)
	\qbezier(52,-1)(52,-3)(51,-3)
	\qbezier(51,-3)(50,-3)(51,0)
	\qbezier(51,0)(55,5)(55,-1)
	\qbezier(55,-1)(55,-3)(54,-3)
	\qbezier(54,-3)(53,-3)(54,0)
	\qbezier(54,0)(58,5)(58,-1)
	\qbezier(58,-1)(58,-3)(57,-3)
	\qbezier(57,-3)(56,-3)(57,0)
	\qbezier(57,0)(59,4)(59,0)
	\put(59,0){\line(1,0){3}}
	% right wall
	\put(62,-5){\line(0,1){10}}
	\put(62,-5){\line(1,0){2}}
	\put(62,5){\line(1,0){2}}
	% displacements
	\put(22,-5){\line(0,1){2}}
	\put(22,-4){\vector(1,0){4}}
	\put(23,-7){\mbox{$x_1$}}
	\put(42,-5){\line(0,1){2}}
	\put(42,-4){\vector(1,0){4}}
	\put(43,-7){\mbox{$x_2$}}
\end{picture}
\end{center}
\caption{Two equal masses attached to each other and fixed walls by
equal springs.}
		\label{fig:twosprings}
\end{figure}


Let $x_1$ denote the displacement of the first mass and $x_2$ the displacement
of the second, and note the displacement of the walls is zero. Each
mass experiences forces from the adjacent springs proportional to the
stretch or compression of the spring.
Ignoring any friction, Newton's law of motion $ma = F$, leads to
\begin{equation}
\begin{array}{rlcclrrl}
m \ddot{x}_1 &= -k(x_1-0)&+k(x_2-x_1)&	       &=& -2k x_1 &+ k x_2 &\text{and}\\
m \ddot{x}_2 &=		 &-k(x_2-x_1)&+k(0-x_2) &=& k x_1 &- 2 k x_2&.\\
\end{array}
\end{equation}
%\begin{equation}
%\begin{split}
%m \ddot{x}_1 & = -k x_1 + k(x_2 - x_1) = -2k x_1 + k x_2 \\
%m \ddot{x}_2 & = -k x_2 - k(x_2 - x_1) = k x_1 - 2 k x_2.
%\end{split}
%\end{equation}
Dividing both sides by $m$ we can write these equations in matrix form
\begin{equation} \label{linearde}
\ddot{\xb} = - A \xb,
\end{equation}
where
\begin{equation}
A   = \left(\begin{array}{cc} 2 \frac{k}{m} & -1 \frac{k}{m} \\
                                                        -1 \frac{k}{m} & 2  \frac{k}{m}
                                       \end{array} \right).
\end{equation}
For this type of equation, the general solution is
\begin{equation}
\xb(t)
= c_1 \vb_1 \sin\left(\sqrt{\lambda_1}\; t + \phi_1\right)
+ c_2 \vb_2 \sin\left(\sqrt{\lambda_2}\; t + \phi_2\right)
\end{equation}
where $\lambda_1$ and $\lambda_2$ are eigenvalues of $A$ with corresponding eigenvectors
$\vb_1$ and $\vb_2$. One can check that this is a solution by substituting it into
the equation (\ref{linearde}).

The eigenvalues of $A$ are the squares of the
frequencies of oscillation. Let's set $m =1$ and $k=1$ in $A$.
We can find the eigenvalues and eigenvectors of $A$ using Matlab:
\begin{lstlisting}[frame=none,numbers=left]
A = [2 -1 ; -1 2]
[v e] = eig(A)
\end{lstlisting}
This should produce a matrix \verb&v& whose columns are the eigenvectors of $A$
and a diagonal matrix $e$ whose entries are the eigenvalues of $A$.
In the first eigenvector, $\vb_1$, the two entries are equal. This represents
the mode of oscillation where the two masses move in sync with each other.
The second eigenvector, $\vb_2$, has the same entries but opposite signs. This represents
the mode where the two masses oscillate in anti-synchronization.
See Figure~\ref{fig:twospringvectors}.
Notice that the frequency for anti-sync motion is $\sqrt{3}$ times that of
synchronous motion.

\begin{figure}
\setlength{\unitlength}{1 mm}
\begin{picture}(64,70)(0,-5)
%% neutral
	%left wall
	\put(0,-5){\line(1,0){2}}
	\put(0,5){\line(1,0){2}}
	\put(2,-5){\line(0,1){10}}
	% spring
	\put(2,0){\line(1,0){3}}
	% weight
	\put(19,0){\line(1,0){6}}
	\put(22,0){\circle*{2}}
	% spring
	%weight
	\put(39,0){\line(1,0){6}}
	\put(42,0){\circle*{2}}
	% spring
	\put(59,0){\line(1,0){3}}
	% right wall
	\put(62,-5){\line(0,1){10}}
	\put(62,-5){\line(1,0){2}}
	\put(62,5){\line(1,0){2}}
%%%%% shift right
	%left wall
	\put(0,10){\line(1,0){2}}
	\put(0,20){\line(1,0){2}}
	\put(2,10){\line(0,1){10}}
	% spring
	\put(2,15){\line(1,0){3}}
	% weight
	\put(23,15){\line(1,0){6}}
	\put(26,15){\circle*{2}}
	% spring
	%weight
	\put(43,15){\line(1,0){6}}
	\put(46,15){\circle*{2}}
	% spring
	\put(59,15){\line(1,0){3}}
	% right wall
	\put(62,10){\line(0,1){10}}
	\put(62,10){\line(1,0){2}}
	\put(62,20){\line(1,0){2}}
%%%%% neutral
	%left wall
	\put(0,25){\line(1,0){2}}
	\put(0,35){\line(1,0){2}}
	\put(2,25){\line(0,1){10}}
	% spring
	\put(2,30){\line(1,0){3}}
	% weight
	\put(19,30){\line(1,0){6}}
	\put(22,30){\circle*{2}}
	% spring
	%weight
	\put(39,30){\line(1,0){6}}
	\put(42,30){\circle*{2}}
	% spring
	\put(59,30){\line(1,0){3}}
	% right wall
	\put(62,25){\line(0,1){10}}
	\put(62,25){\line(1,0){2}}
	\put(62,35){\line(1,0){2}}
%%%%% shift left
	%left wall
	\put(0,40){\line(1,0){2}}
	\put(0,50){\line(1,0){2}}
	\put(2,40){\line(0,1){10}}
	% spring
	\put(2,45){\line(1,0){3}}
	% weight
	\put(15,45){\line(1,0){6}}
	\put(18,45){\circle*{2}}
	% spring
	%weight
	\put(35,45){\line(1,0){6}}
	\put(38,45){\circle*{2}}
	% spring
	\put(59,45){\line(1,0){3}}
	% right wall
	\put(62,40){\line(0,1){10}}
	\put(62,40){\line(1,0){2}}
	\put(62,50){\line(1,0){2}}
%% neutral
	%left wall
	\put(0,55){\line(1,0){2}}
	\put(0,65){\line(1,0){2}}
	\put(2,55){\line(0,1){10}}
	% spring
	\put(2,60){\line(1,0){3}}
	% weight
	\put(19,60){\line(1,0){6}}
	\put(22,60){\circle*{2}}
	% spring
	%weight
	\put(39,60){\line(1,0){6}}
	\put(42,60){\circle*{2}}
	% spring
	\put(59,60){\line(1,0){3}}
	% right wall
	\put(62,55){\line(0,1){10}}
	\put(62,55){\line(1,0){2}}
	\put(62,65){\line(1,0){2}}
\end{picture}
\hfill
\begin{picture}(64,70)(0,-5)
% neutral
	%left wall
	\put(0,-5){\line(1,0){2}}
	\put(0,5){\line(1,0){2}}
	\put(2,-5){\line(0,1){10}}
	% spring
	\put(2,0){\line(1,0){3}}
	% weight
	\put(19,0){\line(1,0){6}}
	\put(22,0){\circle*{2}}
	% spring
	%weight
	\put(39,0){\line(1,0){6}}
	\put(42,0){\circle*{2}}
	% spring
	\put(59,0){\line(1,0){3}}
	% right wall
	\put(62,-5){\line(0,1){10}}
	\put(62,-5){\line(1,0){2}}
	\put(62,5){\line(1,0){2}}
%%%%% together
	%left wall
	\put(0,10){\line(1,0){2}}
	\put(0,20){\line(1,0){2}}
	\put(2,10){\line(0,1){10}}
	% spring
	\put(2,15){\line(1,0){3}}
	% weight
	\put(23,15){\line(1,0){6}}
	\put(26,15){\circle*{2}}
	% spring
	%weight
	\put(35,15){\line(1,0){6}}
	\put(38,15){\circle*{2}}
	% spring
	\put(59,15){\line(1,0){3}}
	% right wall
	\put(62,10){\line(0,1){10}}
	\put(62,10){\line(1,0){2}}
	\put(62,20){\line(1,0){2}}
%%%%% neutral
	%left wall
	\put(0,25){\line(1,0){2}}
	\put(0,35){\line(1,0){2}}
	\put(2,25){\line(0,1){10}}
	% spring
	\put(2,30){\line(1,0){3}}
	% weight
	\put(19,30){\line(1,0){6}}
	\put(22,30){\circle*{2}}
	% spring
	%weight
	\put(39,30){\line(1,0){6}}
	\put(42,30){\circle*{2}}
	% spring
	\put(59,30){\line(1,0){3}}
	% right wall
	\put(62,25){\line(0,1){10}}
	\put(62,25){\line(1,0){2}}
	\put(62,35){\line(1,0){2}}
%%%%% apart
	%left wall
	\put(0,40){\line(1,0){2}}
	\put(0,50){\line(1,0){2}}
	\put(2,40){\line(0,1){10}}
	% spring
	\put(2,45){\line(1,0){3}}
	% weight
	\put(15,45){\line(1,0){6}}
	\put(18,45){\circle*{2}}
	% spring
	%weight
	\put(43,45){\line(1,0){6}}
	\put(46,45){\circle*{2}}
	% spring
	\put(59,45){\line(1,0){3}}
	% right wall
	\put(62,40){\line(0,1){10}}
	\put(62,40){\line(1,0){2}}
	\put(62,50){\line(1,0){2}}
%% neutral
	%left wall
	\put(0,55){\line(1,0){2}}
	\put(0,65){\line(1,0){2}}
	\put(2,55){\line(0,1){10}}
	% spring
	\put(2,60){\line(1,0){3}}
	% weight
	\put(19,60){\line(1,0){6}}
	\put(22,60){\circle*{2}}
	% spring
	%weight
	\put(39,60){\line(1,0){6}}
	\put(42,60){\circle*{2}}
	% spring
	\put(59,60){\line(1,0){3}}
	% right wall
	\put(62,55){\line(0,1){10}}
	\put(62,55){\line(1,0){2}}
	\put(62,65){\line(1,0){2}}
\end{picture}
\caption{
Two vibrational modes of a simple oscillating system. In the left mode
the weights move together and in the right mode they move
opposite. Note that the two modes actually move at different speeds.}
\label{fig:twospringvectors}
\end{figure}


Which of the two modes is the most dangerous for a structure or machine?
It is the one with the {\em lowest frequency} because that mode
can have the largest displacement. Sometimes this mode is called
the {\em fundamental mode}.

We can do the same for three equal masses. With $m=1$, $k=1$ the corresponding matrix $A$
would be
$$
A =  \left(\begin{array}{ccc} 2 & -1 & 0 \\
                              -1& 2 & -1 \\
                              0& -1 & 2
                                       \end{array} \right).
$$
Find the eigenvectors and eigenvalues as above.
There are three different modes. Interpret them from the eigenvectors.



%%%%%%%%%%%%%%%%

\subsection*{Exercises}
\begin{list}
{\arabic{chapter}.\arabic{exercise}}{\usecounter{exercise}}

\item Find the modes and their frequencies for 4 equal masses with $m=3$ kg and equal springs with $k=10$ N/m.
Describe the modes (a sketch will suffice).
%\newpage % so doesn't break problem
\item
Find the modes and their frequencies for  three unequal masses $m_1 = 2$ kg, $m_2 = 3$ kg and $m_3 = 4$ kg
connected by 4 equal springs with $k = 5$ N/m.
How do unequal masses affect the modes? (You must start with the equations of motion to
do this correctly.)

\end{list}


%%%%%%%%%%%%%%%%%%%%%%%%%%%%%%%%%%%%%%%%%%

\chapter{Numerical Methods for Eigenvalues}
\label{EWmeth}

As mentioned above, the eigenvalues and eigenvectors of an $n \times n$ matrix
where $n \ge 4$ must be found numerically instead of by hand.
The numerical methods that are used in practice depend on
the geometric meaning of eigenvalues and eigenvectors which is equation
(\ref{def:eig}). The essence of all these methods is captured
in the Power method, which we now introduce.

%%%%%%%%%%%%%%%%

\subsection*{The Power Method}


In the command window of \textsc{Matlab} enter
\begin{lstlisting}[frame=none,numbers=left]
A = hilb(5)
x = ones(5,1)
x = A*x
el = max(x)
x = x/el
\end{lstlisting}

Compare the new value of \verb&x& with the original.
Repeat the last three lines (you can use the scroll up button).
Compare the newest value of \verb&x& with the previous one and
the original. Notice that there is less change between the second two.
Repeat the last three commands over and over until the values
stop changing. You have completed what is known as the {\em Power Method}.
Now try the command
\begin{lstlisting}[frame=none,numbers=left]
[v e] = eig(A)
\end{lstlisting}
The last entry in \verb&e& should be the final \verb&el& we
computed. The last column in \verb&v& is \verb&x/norm(x)&.
Below we explain why our commands gave this eigenvalue and eigenvector.


For illustration consider a $2 \times 2$ matrix whose
eigenvalues are $1/3$ and $-2$ and whose corresponding eigenvectors
are $\vb_1$ and $\vb_2$.
Let $\xb_0$ be any vector which is a combination of
$\vb_1$ and $\vb_2$, e.g.,
$$
     \xb_0 = \vb_1 + \vb_2.
$$
Now let $\xb_1$ be  $A$ times  $\xb_0$. It follows from (\ref{def:eig}) that
\begin{equation}
\begin{split}
     \xb_1 &= A \vb_1 + A \vb_2 \\
           &= \frac{1}{3} \vb_1 +  -2 \vb_2.
\end{split}
\end{equation}
Thus the $\vb_1$ part is shrunk while the $\vb_2$ is stretched.
If we repeat this process $k$ times then
\begin{equation}
\begin{split}
     \xb_k &= A \xb_{k-1} \\
           &= A^k \xb_0 \\
           &= \left(\frac{1}{3}\right)^k \vb_1 + (-2)^k \vb_2.
\end{split}
\end{equation}
Clearly, $\xb_k$ grows in the direction of $\vb_2$ and
shrinks in the direction of $\vb_1$. This is the principle
of the Power Method, vectors multiplied by $A$
are stretched most in the direction of the eigenvector  whose
eigenvalue has the largest absolute value.

The eigenvalue with the largest absolute value is called the {\em dominant}
eigenvalue. In many applications this quantity will necessarily be positive
for physical reasons. When this is the case, the \textsc{Matlab} code
above will work since \verb&max(v)& will be the element with the
largest absolute value. In applications where the dominant eigenvalue may be
negative, the program must use flow control to determine the correct
number.

Summarizing the Power Method:
\begin{itemize}
\item Repeatedly multiply $\xb$ by $A$ and
  divide by the element with the largest absolute value.
\item The element of largest absolute value converges to largest absolute eigenvalue.
\item The vector converges to the corresponding eigenvector.
\end{itemize}
Note that this logic only works when the eigenvalue largest in
magnitude is real. If the matrix and starting vector are real then the
power method can never give a result with an imaginary part.
Eigenvalues with imaginary part mean the matrix has a rotational
component, so the eigenvector would not settle down either.

Try
\begin{lstlisting}[frame=none,numbers=left]
A = randn(15,15);
e = eig(A)
\end{lstlisting}
You can see that for a random square matrix, many of the eigenvalues are complex.

However, matrices in applications are not just random. They have structure, 
and this can lead to real eigenvalues as seen in the next section.


%%%%%%%%%%%%%%%%%%%%%%%%%

\subsection*{Symmetric, Positive-Definite Matrices}

As noted in the previous paragraph, the power method can fail
if $A$ has complex eigenvalues. One class of matrices that appear often
in applications and for which the eigenvalues are always real are called
the symmetric matrices. A matrix is {\em symmetric} if
$$
     A' = A,
$$
i.e.~$A$ is symmetric with respect to reflections about its diagonal.


Try\begin{lstlisting}[frame=none,numbers=left]
A = rand(5,5)
C = A'*A
e = eig(C)
\end{lstlisting}
You can see that  the eigenvalues of these symmetric matrices are real.


Next we consider an even more specialized class for which the
eigenvalues are not only real, but positive. A symmetric matrix is
called {\em positive definite} if for all vectors $\vb \neq
\mathbf{0}$ the following holds:
$$
     A \vb \cdot \vb > 0.
$$
Geometrically, $A$ does not rotate any vector by more than $\pi/2$.
In summary:
\begin{itemize}
\item If $A$ is symmetric then its eigenvalues are real.
\item If $A$ is symmetric positive definite, then its eigenvalues are positive numbers.
\end{itemize}

Notice that the $B$ matrices in the previous section were symmetric
and the eigenvalues were all real. Notice that the Hilbert and Pascal matrices
are symmetric.

%The matrix $H$ also is positive definite, as are the matrices $B$ in the previous section.


%%%%%%%%%%%%%%%%%%%%%%%%%


\subsection*{The residual of an approximate eigenvector-eigenvalue pair}

If $\vb$ and $\lambda$ are an eigenvector-eigenvalue pair for $A$, then they are supposed to 
satisfy the equations: $A \vb = \lambda \vb$. Thus a scalar residual for approximate
 $\vb$ and $\lambda$ would be
$$
        r  = \| A \vb  - \lambda \vb \|.
$$




%%%%%%%%%%%%%%%%%%%%%%%%%%%%%%%%%%%%%

\subsection*{The Inverse Power Method}

In the application of vibration analysis,  the mode (eigenvector) with
the lowest frequency (eigenvalue) is the most dangerous for
the machine or structure. The Power Method gives us instead
the largest eigenvalue, which is the least important frequency.
In this section we introduce a method, the {\em Inverse
Power Method} which produces exactly what is needed.


The following facts are at the heart of the Inverse Power Method:
\begin{itemize}
\item  If $\lambda$ is an eigenvalue  of $A$ then $1/\lambda$ is an eigenvalue
for $A^{-1}$.
\item The eigenvectors for $A$ and $A^{-1}$ are the same.
\end{itemize}


Thus if we apply the Power Method to $A^{-1}$ we will obtain the largest
absolute eigenvalue of $A^{-1}$, which is exactly the reciprocal of the smallest
absolute eigenvalue of $A$. We will also obtain the corresponding eigenvector, which is
an eigenvector  for both $A^{-1}$ and $A$. Recall that in the application of
vibration mode analysis, the smallest eigenvalue and its eigenvector correspond exactly to
the frequency and mode that we are most interested in, i.e. the one
that can do the most damage.

Here as always, we do not really want to calculate the
inverse of $A$ directly if we can help it. Fortunately, multiplying
$\xb_i$ by $A^{-1}$ to get $\xb_{i+1}$ is equivalent to solving the system
$    A \xb_{i+1} = \xb_i\, $,
which can be done efficiently and accurately. Since iterating this process involves
solving a linear system with the same $A$ but many different right hand
sides, it is a perfect time to use the LU decomposition to save computations.
The following function program does $n$ steps of the Inverse Power Method.
%\newpage
\begin{lstlisting}
function [x e] = myipm(A,n)
    % Performs the inverse power method.
    % Inputs: A -- a square matrix   
    %         n --  number of iterations.
    % Outputs: x -- estimated eigenvector
    %          e -- estimated smallest eigenvalue. 
    [L U P] = lu(A);  % LU decomposition of A with pivoting
    m = size(A,1); % determine the size of A
    x = ones(m,1); % make an initial vector with ones
    for i = 1:n
        px = P*x;  % Apply pivot
        y = L\px;  % solve via LU
        x = U\y;
        % find the maximum entry in absolute value, retaining its sign
        M = max(x);  
        m = min(x);
        if abs(M) >= abs(m)
           el = M;
        else
           el = m;
        end
        x = x/el; % divide by the estimated eigenvalue of the inverse of A
    end
    e = 1/el; % reciprocate to get an eigenvalue of A
end
\end{lstlisting}
Create a $5 \times 5$ Hilbert matrix $H$ and use the program to find its smallest eigenvalue and corresponding eigenvector.
Also do \verb& [V, E] = eig(H) & and compare. 

%%%%%%%%%%%%%%%%%%%%

\subsection*{Exercises}
\begin{list}
{\arabic{chapter}.\arabic{exercise}}{\usecounter{exercise}}

\item For each of the matrices
$$
  A =\left[\begin{array}{cc}
      1&2\\
      3&4
    \end{array}\right]    \quad \text{and} \quad
  B =\left[\begin{array}{ccc}
      -2&1&0\\
      1&-2&1\\
      0&1&-3
    \end{array}\right]
 \,,
$$
 perform two iterations of
  the power method by hand starting with a vector of all ones.
  State the resulting approximations of the eigenvalue and
  eigenvector.
\item 
  \begin{enumerate}
  \item[(a)] Write a well-commented \textsc{Matlab} {\bf function} program
    \verb&mypm.m& that inputs a matrix and a tolerance, applies the
    power method until the scalar residual is less than the
    tolerance, and outputs  the estimated
    eigenvalue and eigenvector, the number of steps, and the scalar residual.
%Use a \verb&while& loop to make it
%    continue until the vector stops changing (up to some \verb&tol&).
%    Put a counter in the loop to count the number of steps.
 \item[(b)] 
   Test your program on the matrices $A$ and $B$ in the previous exercise and check
   that you get the same results as your hand calculations for the first 2 iterations. 
   Turn in your program only.
\end{enumerate}
%Try your program on some
%  random symmetric matrices of various sizes (Hint: use
%  \verb&A=randn(n)& then \verb&A=A+A'&). Check that your program
%  worked by making sure \verb&norm(A*v-e*v)& is small.  Is there a
%  relationship between the size and the number of steps needed or
%  between the size and the eigenvalue? Turn your program and
%  a brief report on the trials.

\end{list}


%%%%%%%%%%%%%%%%%%%%%%%%%%%%%%%%%%%%%%%%%%%%%%%%%%%%%%

\chapter{The QR Method*}

The Power Method and Inverse Power Method each give us
only one eigenvalue-eigenvector pair. While both of these methods can
be modified to give more eigenvalues and eigenvectors, there is
a better method for obtaining all the eigenvalues called
the {\em QR method}. This is the basis of all modern
eigenvalue software, including \textsc{Matlab}, so we summarize
it briefly here.

The QR method uses the fact that any square matrix has a
{\em QR decomposition}. That is, for any $A$ there are
matrices $Q$ and $R$ such the $A = QR$ where $Q$ has the property
$$
     Q^{-1} = Q'
$$
and $R$ is upper triangular. A matrix $Q$ with the property that its
transpose equals its inverse is called an {\em orthogonal} matrix,
because its column vectors are mutually orthogonal.

The QR method consists of iterating following steps:
\begin{itemize}
\item Transform $A$ into a tridiagonal matrix $H$.
\item Decompose $H$ in $QR$.
\item Multiply $Q$ and $R$ together in reverse order to form a new $H$.   
\end{itemize}
The diagonal of $H$ will converge to the eigenvalues.

The details of what makes this method converge are beyond the scope of the this book.
However, we note the following theory behind it for those with more familiarity with
linear algebra. First the Hessian matrix $H$ is
obtained from $A$ by a series of similarity transformation, thus it has the
same eigenvalues as $A$. Secondly, if we denote by $H_0$, $H_1$, $H_2$, $\ldots$, the
sequence of matrices produced by the iteration, then
$$
   H_{i+1} = R_i Q_i = Q_i^{-1} Q_i R_i Q_i = Q'_i H_i Q_i.
$$
Thus each $H_{i+1}$ is a related to $H_i$ by an (orthogonal) similarity transformation and
so they have the same eigenvalues as $A$.

There is a built-in QR decomposition in \textsc{Matlab} which
is called with the command: \verb&[Q R] = qr(A)&.
Thus the following program implements QR method until it
converges:
\newpage
\begin{lstlisting}
function [E,steps] = myqrmethod(A)
    % Computes all the eigenvalues of a matrix using the QR method.
    % Input: A -- square matrix
    % Outputs: E -- vector of eigenvalues 
    %          steps -- the number of iterations it took
    [m n] = size(A);
    if m ~= n
       warning('The input matrix is not square.')
       return
    end
    % Set up initial estimate
    H = hess(A);
    E = diag(H);
    change = 1;
    steps = 0;
    % loop while estimate changes
    while change > 0
        Eold = E;
        % apply QR method
        [Q R] = qr(H);
        H = R*Q;
        E = diag(H);
        % test change
        change = norm(E - Eold);
        steps = steps +1;
    end
end
\end{lstlisting}
As you can see the main steps of the program are very simple.
The really hard calculations are contained in the built-in commands
\verb&hess(A)& and \verb&qr(H)&.

Run this program and compare the results with \textsc{Matlab}'s built in command:
\begin{lstlisting}[frame=none,numbers=left]
format long
format compact
A = hilb(5)
[Eqr,steps] = myqrmethod(A)
Eml = eig(A)
diff = norm(Eml - flipud(Eqr))
\end{lstlisting}

%%%%%%%%%%%%%%%%

\subsection*{Exercises}
\begin{list}
{\arabic{chapter}.\arabic{exercise}}{\usecounter{exercise}}

\item Modify \verb&myqrmethod& to stop after 1000
      iterations. Use the modified program
      on the matrix \verb&A = hilb(n)& with \verb&n& equal to
      10, 50, and 200. Use the norm to compare the results to
      the eigenvalues obtained from \textsc{Matlab}'s built-in program
      \verb&eig&. Turn in your program and a brief
      report on the experiment.


\end{list}




%%%%%%%%%%%%%%%%%%%%%%%%%%%%%%%%%%%%%%%%%%%%%%%%%%%%%%%%%

\chapter{Iterative solution of linear systems*}


%%%%%%%%%%%%%%%%

\subsection*{Newton refinement}

%%%%%%%%%%%%%%%%

\subsection*{Conjugate gradient method}



%%%%%%%%%%%%%%%%%%%%%%%%%%%%%%%%%%%%%%%%%%%%%%

\chapter*{Review of Part II}
\addcontentsline{toc}{chapter}{Review of Part II}
\markboth{REVIEW OF PART II}{}
\subsection*{Methods and Formulas}

%%%%%%%%%%%%%%%%

\subsubsection*{Basic Matrix Theory:}
Identity matrix: $ A I = A$, $I A = A$, and $I \vb = \vb$ \\
Inverse matrix: $ A A^{-1} = I$ and $ A^{-1} A = I$\\
Norm of a matrix: $ \|A\| \equiv \max_{\|\vb\|=1} \|A\vb\|$\\
A matrix may be singular or nonsingular. See Lecture 10.

\subsubsection*{Solving Process:}
Learn the exact Gaussian Elimination algorithm:\\
   Row $j$  $\mapsto$ Row $j$ - (ratio) Row $i$ \\
Gaussian Elimination this way produces LU decomposition\\
Row Pivoting (bigger absolute number on top) \\
Back Substitution


\subsubsection*{Condition number:}
$$
      \textrm{cond}(A) \equiv \max \left( \frac{ \|\delta \xb\|/\|\xb\|}
         { \frac{\|\delta A\|}{\|A\|} + \frac{\|\delta \bb\|}{\|\bb\|} } \right)
          = \max \left( \frac{\text{Relative error of output}}
                             {\text{Relative error of inputs}} \right).
$$
A big condition number is bad; in engineering it usually results from poor design.


\subsubsection*{LU factorization:}
The LU factorization is a by-product of Gaussian Elimination (if done with
the correct algorithm).
$$
   PA = LU.
$$
Solving steps:\\
Multiply by P: $\mathbf{d}= P \bb$ \\
Forwardsolve: $  L \yb = \mathbf{d} $ \\
Backsolve: $  U \xb = \yb $


\subsubsection*{Eigenvalues and eigenvectors:}
A  nonzero vector $\vb$ is an eigenvector
and a number $\lambda$ is its eigenvalue if
$$
   A \vb = \lambda \vb.
$$
Characteristic equation: $ \det(A - \lambda I) = 0 $  \\
Equation of the eigenvector: $(A - \lambda I) \vb = \mathbf{0}$\\
Residual for an approximate eigenvector-eigenvalue pair: $ r  = \| A \vb  - \lambda \vb \|$


\subsubsection*{Complex eigenvalues:}
Occur in conjugate pairs: $ \lambda_{1,2} = \alpha \pm i \beta $\\
and  eigenvectors must also come
in conjugate pairs: $\wb = \ub \pm i \vb$.

\subsubsection*{Vibrational modes:}
Eigenvalues are frequencies squared. Eigenvectors represent modes.

\subsubsection*{Power Method:}
- Repeatedly multiply $\xb$ by $A$ and
  divide by the element with the largest absolute value.\\
- The element of largest absolute value converges to largest absolute eigenvalue.\\
- The vector converges to the corresponding eigenvector.\\
- Convergence assured for a real symmetric matrix, but not for an arbitrary
matrix, which may not have real eigenvalues at all.


\subsubsection*{Inverse Power Method:}
- Apply power method to $A^{-1}$.\\
- Use solving rather than the inverse.\\
- If $\lambda$ is an eigenvalue  of $A$ then $1/\lambda$ is an eigenvalue  for $A^{-1}$.\\
- The eigenvectors for $A$ and $A^{-1}$ are the same.

\subsubsection*{Symmetric and Positive definite:}
- Symmetric: $A = A'$.\\
- If $A$ is symmetric its eigenvalues are real.\\
- Positive definite: $A\xb \cdot \xb >0$.\\
- If $A$ is positive definite, then its eigenvalues are positive.


\subsubsection*{QR method (Not covered in MATH 3600 at Ohio):}
- Transform $A$ into $H$ the Hessian form of $A$.\\
- Decompose $H$ in $QR$.\\
- Multiply $Q$ and $R$ together in reverse order to form a new $H$.\\
- Repeat\\
- The diagonal of $H$ will converge to the eigenvalues of $A$.


\subsection*{Matlab}

\subsubsection*{Matrix arithmetic:}

% \begin{lstlisting}[frame=none,numbers=left]
% A = [ 1  3 -2 5 ;  -1  -1 5 4 ; 0 1 -9  0]  % Manually enter a matrix 
% u = [ 1  2  3  4]' 
% A*u  
% B = [3 2 1; 7 6 5; 4 3 2] 
% B*A      % multiply $B$ times $A$
% 2*A      % multiply a matrix by a scalar
% A + A    % add matrices 
% A + 3    % add 3 to every entry of a matrix 
% B.*B     % component-wise multiplication
% B.^3     % component-wise exponentiation
%\end{lstlisting}
\verb&A = [ 1  3 -2 5 ;  -1  -1 5 4 ; 0 1 -9  0]& \dotfill Manually enter a matrix. \\
\verb&u = [ 1  2  3  4]'&\\
\verb&A*u &\\
\verb&B = [3 2 1; 7 6 5; 4 3 2]&\\
\verb&B*A&     \dotfill multiply $B$ times $A$.\\
\verb&2*A&     \dotfill multiply a matrix by a scalar.\\
\verb&A + A&   \dotfill add matrices. \\
\verb&A + 3&   \dotfill add 3 to every entry of a matrix. \\
\verb&B.*B&  \dotfill component-wise multiplication. \\
\verb&B.^3&  \dotfill component-wise exponentiation.

\subsubsection*{Special matrices:}

% \begin{lstlisting}[frame=none,numbers=left]
% I = eye(3)        % identity matrix
% D = ones(5,5) 
% O = zeros(10,10) 
% C = rand(5,5)     % random matrix with uniform distribution in [0,1]
% C = randn(5,5)    % random matrix with normal distribution.
% hilb(6) 
% pascal(5) 
% \end{lstlisting}
\verb&I = eye(3)& \dotfill identity matrix\\
\verb&D = ones(5,5)&\\
\verb&O = zeros(10,10)&\\
\verb&C = rand(5,5)& \dotfill random matrix with uniform
         distribution in $[0,1]$.\\
\verb&C = randn(5,5)& \dotfill random matrix with normal distribution.\\
\verb&hilb(6)&\\
\verb&pascal(5)&

\subsubsection*{General matrix commands:}

% \begin{lstlisting}[frame=none,numbers=left]
% size(C)   % gives the dimensions (m x n) of A
% norm(C)   % gives the norm of the matrix
% det(C)    % the determinant of the matrix
% max(C)    % the maximum of each row
% min(C)    % the minimum in each row
% sum(C)    % sums each row
% mean(C)   % the average of each row
% diag(C)   % just the diagonal elements 
% inv(C)    % inverse of the matrix
% C'        % transpose of the matrix
% \end{lstlisting}
\verb&size(C)& \dotfill gives the dimensions ($m \times n$) of $A$.\\
\verb&norm(C)& \dotfill gives the norm of the matrix.\\
\verb&det(C)&   \dotfill the determinant of the matrix.\\
\verb&max(C)&   \dotfill the maximum of each row.\\
\verb&min(C)&   \dotfill the minimum in each row.\\
\verb&sum(C)&   \dotfill sums each row.\\
\verb&mean(C)&   \dotfill the average of each row.\\
\verb&diag(C)&   \dotfill just the diagonal elements. \\
\verb&inv(C)&    \dotfill inverse of the matrix.\\
\verb&C'&    \dotfill transpose of the matrix.

\subsubsection*{Matrix decompositions:}

%\begin{lstlisting}[frame=none,numbers=left]
%[L U P] = lu(C) 
%[Q R] = qr(C) 
%H = hess(C)     % transform into a Hessian tri-diagonal matrix
%\end{lstlisting}
\verb&[L U P] = lu(C)&\\
\verb&[Q R] = qr(C)&\\
\verb&H = hess(C)& \dotfill transform into a Hessian tri-diagonal matrix, which has the same eigenvalues as $A$.

%[U S V] = svd(C)   % singular value decomposition (important, but we did not use it)
%R = chol(A)   % Cholesky decomposition of a symmetric, positive definite matrix, A = R' R.
%[U T] = schur(C)   % Schur Decomposition A=U'T U.
%H = hess(C)     % transform into a Hessian tri-diagonal matrix, which has the same eigenvalues as A.




%%%%%%%%%%%%%%%%%%%%%%%%%%%%%%%%%%%%%%%%%%%%%%%%%%%%%%%%%%%%%%%%%%%%%%%%%
%%%%%%%%%%%%%%%%%%%%%%%%%%%%%%%%%%%%%%%%%%%%%%%%%%%%%%%%%%%%%%%%%%%%%%%%%
